% Addendum to "The finite-degree profile of essential dimension of a Z/p^n-torsor
% in characteristic p": Section 10 (characteristic 0), Section 11 (supplements at
% (2,2) and (2,3)) and Appendices A-F.  Compiled standalone with page/section
% counters continuing the main paper, and appended to the original PDF.
% To integrate into the source: \input the three body files after Section 9
% (placement comments at the top of body_sec11.tex say where each subsection
% belongs) and append the bibliography items [17]-[30].
\documentclass[11pt]{amsart}
\usepackage[margin=1.15in]{geometry}
\usepackage{amsmath,amssymb,amsthm}
\usepackage{array,booktabs}
\usepackage[hidelinks]{hyperref}

\newtheorem{theorem}{Theorem}[section]
\newtheorem{proposition}[theorem]{Proposition}
\newtheorem{lemma}[theorem]{Lemma}
\newtheorem{corollary}[theorem]{Corollary}
\theoremstyle{definition}
\newtheorem{definition}[theorem]{Definition}
\newtheorem{example}[theorem]{Example}
\theoremstyle{remark}
\newtheorem{remark}[theorem]{Remark}
\newtheorem{question}[theorem]{Question}

\newcommand{\ed}{\operatorname{ed}}
\newcommand{\edd}[1]{\ed^{[#1]}}
\newcommand{\tr}{\operatorname{trdeg}}
\newcommand{\im}{\operatorname{im}}
\newcommand{\Ind}{\operatorname{Ind}}
\newcommand{\gen}{\mathrm{gen}}
\newcommand{\Z}{\mathbb Z}
\newcommand{\Q}{\mathbb Q}
\newcommand{\A}{\mathbb A}
\newcommand{\Gam}{\Gamma}
\newcommand{\stag}[1]{\hfill{\normalfont[#1]}}
\newcommand{\res}{\operatorname{res}}
\newcommand{\cores}{\operatorname{cores}}

\setcounter{section}{9}
\setcounter{page}{17}

\begin{document}
\pagestyle{myheadings}
\markboth{THE FINITE-DEGREE PROFILE OF ESSENTIAL DIMENSION}{THE FINITE-DEGREE PROFILE OF ESSENTIAL DIMENSION}

\section{The profile in arbitrary characteristic, and the characteristic-zero picture}

\noindent\textit{Addendum (27 September 2026).} Sections 1--9 treat $G=\Z/p^n$ in
characteristic $p$. This section records what the definition of the profile gives for
an arbitrary finite group $G$ over an arbitrary field $k$, and what it looks like in
characteristic $0$ (more generally, in characteristic prime to $|G|$), where the
mechanisms are entirely different: instead of Artin--Schreier layers, the profile is
governed by Sylow divisibility (Theorem~\ref{thm:div}) and by the roots-of-unity
defect $\ed_k(G)-\ed_{k(\mu_e)}(G)$ (Corollary~\ref{cor:rou}). Statements carry the
tags of Section~1; in addition, \S\ref{ss:prov} says explicitly which statements are
rewrites of results in the literature and which are new. As in the rest of the note,
literature is quoted without re-verification and the corresponding statements are
tagged [cited].

\subsection{Conventions}\label{ss:conv}
$k$ is any field, $G$ a finite group (constant group scheme over $k$), $F\supseteq k$
a field, $\tau\colon T\to\operatorname{Spec}F$ a $G$-torsor. The classifying
homomorphism $\varphi_\tau\colon\Gam_F\to G$ is defined up to conjugation, its image
$H=\im\varphi_\tau$ up to conjugacy, and $T\cong\coprod_{G/H}\operatorname{Spec}E$
with $E/F$ Galois with group $H$ (the \emph{splitting field} of $\tau$). We use, in
every characteristic: (B1) $\ed_k(\tau\otimes_FF')\le\ed_k(\tau)$; (B2)
$\ed_k(\tau)\ge\ed_{k'}(\tau)$ for $k\subseteq k'\subseteq F$; (B3) $\ed_k(\tau)=0$
iff $\tau$ descends to the algebraic closure of $k$ in $F$ (for $k$ algebraically
closed: iff $\tau$ is split; for general $k$ this is \emph{not} ``descends to $k$'':
$k=\Q$, $F=\Q(i)(t)$, $\tau=F(\sqrt{1+i})$ has $\ed_\Q=0$ but is not
$F(\sqrt a)$ with $a\in\Q$, since $N_{\Q(i)/\Q}(1+i)=2$ is not a square); (B4)
Lemma~4.1(i), purely transcendental base change does not change $\ed_k$. For a prime
$q$, $\ed_k(\tau;q)=\min\{\ed_k(\tau\otimes F'):F'/F\text{ finite separable},\
q\nmid[F':F]\}$ and $\ed_k(G;q)=\sup_\tau\ed_k(\tau;q)$ \cite[(2.5)]{RS}.
The generic torsor $\tau_\gen\colon\operatorname{Spec}L\to\operatorname{Spec}K$,
$L=k(V)$, $K=L^G$, is taken for a faithful representation $V$ (or any faithful
generically free $G$-variety with $k(V)/k$ regular); $\ed_k(\tau_\gen)=\ed_k(G)$
\cite{BF,Me}. The profile $\edd d_k(\tau)$ and the jump degrees $d_j(\tau)$ are as in
Definition~3.1.

\subsection{Formal properties, valid in every characteristic}

\begin{proposition}[shape]\label{prop:shape}
Let $\tau$ be a $G$-torsor over $F\supseteq k$.
\begin{enumerate}
\item[(a)] $d\mapsto\edd d_k(\tau)$ is non-increasing with values in
$\{0,\dots,\ed_k(\tau)\}$, $\edd1_k(\tau)=\ed_k(\tau)$, and
$d_{\ed(\tau)}=1\le d_{\ed(\tau)-1}\le\dots\le d_0\le|\im\varphi_\tau|$.
\item[(b)] If $[F':F]=e$ and $d\ge1$ then $\edd{de}_k(\tau)\le\edd d_k(\tau\otimes
F')\le\edd d_k(\tau)$; hence $d_j(\tau)\le[F':F]\,d_j(\tau\otimes F')$ for every
finite $F'/F$ and every $j$.
\item[(c)] $\edd d_k(\tau)=\min\{j:d_j(\tau)\le d\}$: the jump degrees determine the
profile.
\item[(d)] If $k\subseteq k'\subseteq F$ then $\edd d_{k'}(\tau)\le\edd d_k(\tau)$
for all $d$.
\item[(e)] $\edd d_k(\tau)=0$ iff some $F'/F$ of degree $\le d$ makes $\tau\otimes
F'$ descend to the algebraic closure of $k$ in $F'$; for $k$ algebraically closed, iff
$\tau$ is split by an extension of degree $\le d$.
\end{enumerate}
\stag{proved}
\end{proposition}

\begin{proof}
(a) Over the splitting field $E$, $[E:F]=|H|$, the torsor is split, so
$d_0\le|H|$; the rest is (B1). (b) is Proposition~3.2(b) verbatim; if
$[F'':F']=d_j(\tau\otimes F')$ realises $\ed(\tau\otimes F'')\le j$ then
$[F'':F]=[F':F]\,d_j(\tau\otimes F')$. (c) is a restatement, (d) is (B2) applied to
each $\tau\otimes F'$, (e) is (B3).
\end{proof}

\begin{proposition}[the last jump degree]\label{prop:d0}
Let $H=\im\varphi_\tau$ and $E$ the splitting field of $\tau$.
\begin{enumerate}
\item[(a)] $d_0(\tau)\le\min\{[F':F]:\tau\otimes F'\text{ split}\}=|H|$, with
equality if $k$ is algebraically closed.
\item[(b)] Let $k$ be perfect and $E/k$ regular. If $[F':F]<\infty$ and
$\tau\otimes F'$ descends to the algebraic closure $k_1$ of $k$ in $F'$, then
$[F':F]\ge|H|\cdot[k_1:k]$. In particular $d_0(\tau)=|H|$.
\item[(c)] For $\tau_\gen$ over a perfect $k$: $d_0(\tau_\gen)=|G|$, and
$\ed_k(\tau_\gen\otimes F')=0$ forces $[F':K]\ge|G|\,[k_1:k]$.
\end{enumerate}
\stag{proved}
\end{proposition}

\begin{proof}
(a) $T\otimes_FF'$ is split iff $E$ embeds $F$-linearly into $F'$, which forces
$[F':F]\ge|H|$; $F'=E$ works; for $k$ algebraically closed use (B3).
(b) Let $\tau\otimes F'\cong\tau_0\otimes_{k_1}F'$ with $\tau_0$ over $k_1$, of
splitting field $l/k_1$ Galois with group $H_0$. Comparing total spaces inside an
algebraic closure of $F'$ we may arrange $EF'=lF'$. Since $l/k$ is finite separable
and $E/k$ regular, $l\otimes_kE$ is a field and $[lE:E]=[l:k]=|H_0|[k_1:k]$; as
$lE\subseteq EF'$,
\[
|H_0|[k_1:k]\,|H|=[lE:E][E:F]\le[EF':F]=[lF':F'][F':F]\le|H_0|\,[F':F].
\]
(c) $E=L$ is regular over $k$ and $H=G$.
\end{proof}

\begin{remark}
Without regularity $d_0<|H|$ is possible ($k=\Q$, $\tau=\Q(t)(\sqrt2)$ is constant,
$d_0=1$). Proposition~\ref{prop:d0}(c) is the general form of Proposition~3.2(c). For
$k$ algebraically closed, $\edd d(\tau)\ge1$ for $d<|H|$: the last slot of the profile
is the only one understood in general.
\end{remark}

\begin{proposition}[base-field extension bound]\label{prop:bfe}
Let $k'/k$ be finite of degree $d$. Then $\edd d_k(\tau)\le\ed_{k'}(G)$. More
precisely, every field factor $F'$ of $F\otimes_kk'$ has $[F':F]\le d$ and
$\ed_k(\tau\otimes F')\le\ed_{k'}(\tau\otimes F')\le\ed_{k'}(G)$.
\stag{proved}
\end{proposition}

\begin{proof}
$F'\supseteq k'$ and $[F':F]\le d$. A descent $\tau\otimes F'\cong\tau_0\otimes_{F_0}
F'$ with $k'\subseteq F_0\subseteq F'$ and $\tr_{k'}F_0\le\ed_{k'}(G)$ is also a
descent over $k$, and $\tr_kF_0=\tr_{k'}F_0$ since $k'/k$ is algebraic.
\end{proof}

\begin{corollary}[roots-of-unity defect]\label{cor:rou}
Let $e$ be the exponent of $G$, $\operatorname{char}k\nmid e$, and
$d_e=[k(\mu_e):k]$. For every $G$-torsor $\tau$ over every $F\supseteq k$,
\[
\edd{[F(\mu_e):F]}_k(\tau)\le\edd{d_e}_k(\tau)\le\ed_{k(\mu_e)}(G).
\]
In particular: (i) if $G$ is abelian, $\edd{d_e}_k(\tau)\le\operatorname{rank}G$
\cite{BR}; (ii) if $G$ is a $p$-group and $\operatorname{char}k\ne p$,
$\edd{p-1}_k(\tau)\le\edd{[k(\mu_p):k]}_k(\tau)\le\ed_{k(\mu_p)}(G)$, the minimal
dimension of a faithful representation of $G$ over $k(\mu_p)$ \cite{KM}.
\stag{proved; the values are cited}
\end{corollary}

\begin{remark}
For $\tau_\gen$ and $F'=K(\mu_e)$ the torsor $\tau_\gen\otimes F'$ is the generic
torsor over $k(\mu_e)$, so $\ed_{k(\mu_e)}(\tau_\gen\otimes K(\mu_e))=\ed_{k(\mu_e)}(G)$
exactly, but whether $\ed_k$ of this torsor equals $\ed_{k(\mu_e)}(G)$ is not known to
us in general. [open]
\end{remark}

\subsection{Sylow divisibility of the jump degrees}

\begin{lemma}[$q$-closure]\label{lem:qcl}
Let $q$ be a prime, $S\le\Gam_F$ a Sylow pro-$q$ subgroup and $F^{(q)}=(F^{\rm sep})^S$.
Then $F^{(q)}$ is a union of finite subextensions of degree prime to $q$, every finite
separable extension of $F^{(q)}$ has $q$-power degree, and
$\ed_k(\tau;q)=\ed_k(\tau\otimes F^{(q)})$ \cite[(2.5)]{RS}. If $F'/F$ is finite
separable with $q\nmid[F':F]$, then there is an $F$-embedding $F'\hookrightarrow
F^{(q)}$, and hence $\ed_k(\tau\otimes F')\ge\ed_k(\tau;q)$.
\stag{cited; last statement proved}
\end{lemma}

\begin{proof}
$F'\otimes_FF^{(q)}$ is a product of finite separable extensions of $F^{(q)}$, each of
$q$-power degree, whose degrees sum to $[F':F]$; if none had degree $1$ the sum would
be divisible by $q$. So some factor is $F^{(q)}$, giving the embedding, and (B1)
gives the inequality. (For a \emph{given} embedding $F'\subseteq F^{\rm sep}$ the
compositum $F'F^{(q)}$ need not be $F^{(q)}$: $F=\Q$, $q=2$, $F'=\Q(\sqrt[3]2)$ has
$F'\otimes_\Q\Q^{(2)}\cong\Q^{(2)}\times\Q^{(2)}(\sqrt{-3})$.)
\end{proof}

\begin{lemma}[purely inseparable invariance]\label{lem:pi}
Let $\operatorname{char}k=p$, $k$ perfect, $G$ any finite group, $F'/F$ finite purely
inseparable. Then $\ed_k(\tau\otimes F')=\ed_k(\tau)$; consequently
$\edd d_k(\tau)$ is the minimum over \emph{separable} $F'/F$ of degree $\le d$.
\stag{proved}
\end{lemma}

\begin{proof}
As in Proposition~3.6, with $\sigma\colon F^{1/p^m}\to F$, $x\mapsto x^{p^m}$; the
$p^m$-power map extends to $F^{1/p^m}F^{\rm sep}\to F^{\rm sep}$ and commutes with all
automorphisms, so the induced isomorphism $\Gam_{F^{1/p^m}}\to\Gam_F$ is the canonical
one and the classifying homomorphism of $\sigma_*(\tau\otimes F^{1/p^m})$ is
$\varphi_\tau$. This recovers Proposition~3.6 without Witt vectors.
\end{proof}

\begin{proposition}\label{prop:genq}
For every prime $q$, $\ed_k(\tau_\gen;q)=\ed_k(G;q)$.
\stag{cited: \cite{RS,Me}; sketch below}
\end{proposition}

\begin{proof}[Sketch]
``$\le$'' is the definition. For ``$\ge$'' let $e=\ed_k(\tau_\gen;q)$, realised by a
prime-to-$q$ extension $K'/K$ and a descent to $K_0$, $\tr_kK_0\le e$; spread out to a
$G$-torsor $Z\to B$, $\dim B\le e$, and to a $G$-equivariant rational map
$W'\dashrightarrow Z$ from the normalisation $W'$ of $V$ in $L\otimes_KK'$, finite of
degree $[K':K]$ over $V$ and \'etale over a dense $G$-stable open $V_1\subseteq V$.
Given any torsor $\tau$ over an infinite $F$, the twist ${}^\tau V$ is an affine space
over $F$, so a free $F$-point in ${}^\tau V_1$ exists and gives an equivariant
$f\colon T\to V_1$; the pullback $T\times_VW'\to T$ is finite \'etale of degree
$[K':K]$, hence $\cong T\times_FA$ for a finite \'etale $F$-algebra $A$ of that degree.
Some field factor $F'$ of $A$ has $q\nmid[F':F]$ and $T_{F'}\to W'\dashrightarrow Z$ is
equivariant into a free $G$-variety, so $\ed_k(\tau\otimes F')\le\dim B\le e$ by
Lemma~2.3 (valid in any characteristic, Theorem~\ref{thm:geom}).
\end{proof}

\begin{theorem}[Sylow divisibility]\label{thm:div}
Let $\operatorname{char}k=0$, or $k$ perfect (then use Lemma~\ref{lem:pi}). Let
$\tau$ be a $G$-torsor over $F$.
\begin{enumerate}
\item[(a)] For every prime $q$ and every $j<\ed_k(\tau;q)$, $q$ divides $d_j(\tau)$:
every $F'/F$ with $\ed_k(\tau\otimes F')\le j$ has $q\mid[F':F]$.
\item[(b)] Every prime $q$ with $\ed_k(G;q)>j$ divides $d_j(\tau_\gen)$; hence
$d_j(\tau_\gen)\ge\prod_{q:\,\ed_k(G;q)>j}q$.
\item[(c)] For every $d\ge1$,
$\displaystyle\edd d_k(\tau_\gen)\ \ge\ \min_{1\le d'\le d}\ \max_{q\nmid d'}\ed_k(G;q)
\ \ge\ \max_{q>d}\ed_k(G;q)$ ($q$ prime).
\item[(d)] The prime-to-$q$ profile $\edd d_k{}'(\tau)=\min\{\ed_k(\tau\otimes F'):
[F':F]\le d,\ q\nmid[F':F]\}$ is non-increasing with terminal value $\ed_k(\tau;q)$,
reached at a finite prime-to-$q$ degree; for $\tau_\gen$ the terminal value is
$\ed_k(G;q)$.
\end{enumerate}
\stag{proved, granting Proposition~\ref{prop:genq}}
\end{theorem}

\begin{proof}
(a) is Lemma~\ref{lem:qcl}; (b) adds Proposition~\ref{prop:genq}; (c) a witness for
$\edd d$ has some degree $d'\le d$ and (a) applies to every $q\nmid d'$; (d) is
Lemma~\ref{lem:qcl} together with the finiteness argument of Proposition~3.7(a).
\end{proof}

\begin{remark}
(i) In the language of Theorem~4.2 and Remark~4.5, (a) is Reichstein--Scavia's
Theorem~10.1 applied to the correspondence of degree prime to $q$ furnished by a free
closed point; the proof above via the $q$-closure is the same mechanism.
(ii) Consistency with Sections 3--9: for $G=\Z/p^n$ in characteristic $p$,
$\ed_k(G;p)\le1$ and $\ed_k(G;q)=0$ for $q\ne p$, so Theorem~\ref{thm:div} only says
$p\mid d_0$, consistent with $d_0=p^n$; it says nothing about $d_1,\dots,d_{n-1}$,
which is the point of this note. In characteristic $0$ the constraint is much stronger.
\end{remark}

\begin{proposition}[$q$-groups]\label{prop:qgp}
Let $G$ be a $q$-group of order $q^n$, $\operatorname{char}k\ne q$ ($k$ perfect if
positive), $c=[k(\mu_q):k]$. Then:
(a) $\edd d_k(\tau_\gen)=\ed_k(G;q)$ for $c\le d\le q-1$, and
$\ed_k(G;q)\le\edd d_k(\tau_\gen)\le\ed_k(G)$ for $d<c$;
(b) $\ed_k(G;q)=\ed_{k(\mu_q)}(G)$ is the minimal dimension of a faithful
$k(\mu_q)$-representation \cite{KM};
(c) $q\mid d_j(\tau_\gen)$ for every $j<\ed_k(G;q)$, and $d_0=q^n$;
(d) if $\mu_q\subset k$: $\edd d(\tau_\gen)=\ed_k(G)$ for $d<q$ and all steps of the
profile occur at multiples of $q$, i.e.\ $\edd d(\tau_\gen)=\edd{q\lfloor d/q\rfloor}
(\tau_\gen)$ for $d\ge q$.
\stag{proved, modulo \cite{KM} and $\ed(\cdot;q)$ being invariant under prime-to-$q$
extensions of $k$}
\end{proposition}

\begin{proof}
Lower bounds: every degree $<q$ is prime to $q$ (Theorem~\ref{thm:div}). Upper bound
for $d\ge c$: Corollary~\ref{cor:rou} and $\ed_{k(\mu_q)}(G)=\ed_{k(\mu_q)}(G;q)=
\ed_k(G;q)$ \cite{KM,Me}. (c) is Theorem~\ref{thm:div}(b) and
Proposition~\ref{prop:d0}(c); (d) is (a)+(c) with $c=1$.
\end{proof}

Compare Proposition~3.7: in characteristic $p$ the prime-to-$p$ tail of a $p$-group
collapses to $\le1$; in characteristic $\ne q$ the prime-to-$q$ tail of a $q$-group is
\emph{constant}, equal to $\ed_k(G;q)$.

\subsection{Geometric meaning and induced torsors in every characteristic}

\begin{theorem}\label{thm:geom}
Theorem~4.2 and Corollary~4.4 hold for every finite group $G$, every field $k$ and
every characteristic: $\edd d_k(\tau)\le m$ iff some quasi-projective generically free
primitive $G$-variety $Z$ of dimension $\le m$ has a closed point of degree $\le d$ in
the free locus of ${}^\tau Z$, iff there are such $Z$, $F'/F$ of degree $\le d$ and a
$G$-equivariant morphism $T_{F'}\to Z^{\rm free}$. For $\tau_\gen$ this is a
$G$-correspondence $V\rightsquigarrow Z$ of degree $\le d$ into $Z^{\rm free}$.
\stag{proved}
\end{theorem}

\begin{proof}
The proof of Theorem~4.2 uses only: the existence of the twist ${}^\tau Z=(T\times_kZ)/G$
as a quasi-projective $F$-scheme compatible with base change and open immersions
\cite[\S3]{DR}; the identification of $F$-points of ${}^\tau Z$ with equivariant maps
$T\to Z_F$ (elementary for finite $G$); the fact that $Z^{\rm free}\to Z^{\rm free}/G$
is a $G$-torsor (SGA~1, V, finite group acting freely); and spreading out. None of
these depends on the characteristic or on $G$ being a cyclic $p$-group.
\end{proof}

\begin{lemma}[induced torsors]\label{lem:ind}
For finite groups $H\le G$ and an $H$-torsor $\tau_H$ over $F$,
$\ed_k(\Ind_H^G\tau_H)=\ed_k(\tau_H)$.
\stag{proved}
\end{lemma}

\begin{proof}
``$\le$'': descend and induce. ``$\ge$'': reduce to $T_H$ a field (if
$T_H=\Ind_{H'}^HT'$ with $T'$ a field, then $\Ind_H^G\tau_H=\Ind_{H'}^G\tau'$ and
$\ed(\tau_H)\le\ed(\tau')$). Then the argument of Proposition~3.9 goes through for
non-abelian $G$: a descent $T_0\to\operatorname{Spec}F_0$ has a connected component
$T_0'$ with stabiliser $H_0$; the components of $T_0'\otimes_{F_0}F$ are among the
$gT_H$ and are permuted transitively by $H_0$, so after conjugating $H$ we may assume
$T_H\subseteq T_0'\otimes_{F_0}F$ and $H\le H_0$. With $F_1=(T_0')^H$ and
$F_1\otimes_{F_0}F=\prod F_i$, the factor $T_0'\otimes_{F_1}F_{i_0}$ containing $T_H$
is an $H$-torsor over $F_{i_0}$ on whose components $H$ acts transitively; being
$H$-stable, $T_H$ is that factor, and $F_{i_0}=T_H^H=F$. So $\tau_H$ descends to
$F_1$, $\tr_kF_1=\tr_kF_0$.
\end{proof}

\begin{corollary}[subgroup profile]\label{cor:sub}
For every subgroup $H'\le G$: $\edd{[G:H']}_k(\tau_\gen)\le\ed_k(\tau_\gen\otimes
L^{H'})=\ed_k(H')$. Hence $d_j(\tau_\gen)\le\min\{[G:H']:\ed_k(H')\le j\}$.
\stag{proved}
\end{corollary}

\begin{proof}
$\tau_\gen\otimes L^{H'}=\Ind_{H'}^G(L/L^{H'})$, and $L/L^{H'}$ is the generic
$H'$-torsor of $V|_{H'}$ over $k$; apply Lemma~\ref{lem:ind}.
\end{proof}

This is the general form of Proposition~3.9 and of the root-field reductions of
Example~\ref{ex:Sn}.

\subsection{Examples in characteristic zero}

\begin{example}[$\Z/n$]\label{ex:cyc}
Let $\mu_n\subset k$, $\operatorname{char}k\nmid n$. The generic torsor
$K(\sqrt[n]a)/K$, $K=k(a)$, has $\ed=1$ and $d_0=n$ (Proposition~\ref{prop:d0}), so
its profile is $(1,\dots,1,0)$. \stag{proved}
\end{example}

\begin{theorem}[$(\Z/2)^n$]\label{thm:ea}
Let $k$ be algebraically closed of characteristic $\ne2$, $G=(\Z/2)^n$,
$F=K=k(a_1,\dots,a_n)$, $\tau_\gen=(a_1,\dots,a_n)\in H^1(F,\Z/2)^n$.
\begin{enumerate}
\item[(a)] $\edd{2^j}(\tau_\gen)\le n-j$ for $0\le j\le n$, and $d_0=2^n$.
\item[(b)] Every $F'/F$ of odd degree has $\ed(\tau_\gen\otimes F')=n$; all $d_j$,
$j<n$, are even.
\item[(c)] For every quadratic $F'/F$, $\ed(\tau_\gen\otimes F')\ge n-1$. Hence
$\edd2=\edd3=n-1$, $d_{n-1}=2$, and $d_{n-2}=4$ for $n\ge2$.
\item[(d)] $\edd d=1$ for $2^{n-1}\le d<2^n$. Complete profiles: $(1,0)$ for $n=1$;
$(2,1,1,0)$ for $n=2$; $(3,2,2,1,1,1,1,0)$ for $n=3$.
\end{enumerate}
\stag{proved; inputs cited}
\end{theorem}

\begin{proof}
(a) Over $F(\sqrt{a_1},\dots,\sqrt{a_j})$ the class is $(1,\dots,1,a_{j+1},\dots,a_n)$.
(b) $\ed_k(G;2)=\ed_k(G)=n$ \cite{BR,KM}; apply Theorem~\ref{thm:div}. (d) from (a),
(b), (c) and Proposition~\ref{prop:d0}.

(c) Write $(a_I)=\bigcup_{i\in I}(a_i)\in H^{|I|}(F,\Z/2)$. Suppose
$\ed(\tau_\gen\otimes F')\le n-2$ for $F'=F(\sqrt f)$. The class descends to $F_0$ with
$\tr_kF_0\le n-2$, so $H^{n-1}(F_0,\Z/2)=0$ \cite[II.4.2]{Se} and $(a_I)_{F'}=0$ for
all $|I|=n-1$. By the Arason exact sequence $H^{r-1}(F)\xrightarrow{(f)\cup}H^r(F)
\xrightarrow{\res}H^r(F(\sqrt f))$ \cite{Ar}, $(a_{[n]\smallsetminus i})\in(f)\cup
H^{n-2}(F)$ for all $i$. We show by induction on $m$ that for $F_m=k(a_1,\dots,a_m)$
there is no $g\in F_m^\times$ with $(a_{[m]\smallsetminus i})\in(g)\cup H^{m-2}(F_m)$
for all $i\in[m]$. For $m=1$: $(a_\emptyset)=1\ne0$ in $H^0$. For $m\ge2$ pass to
$\hat F=F_{m-1}((a_m))$, where $H^r(\hat F)=H^r(F_{m-1})\oplus H^{r-1}(F_{m-1})\cup(a_m)$
and units have $(u)=(\bar u)$ \cite[II.4.3]{Se}; the classes $(a_I)$ restrict
injectively (their residues are nonzero by induction on the residue maps). Write
$g=a_m^\varepsilon u$ and the given $x_i=y_i+z_i\cup(a_m)$. If $\varepsilon=1$, take
$i=m$: comparing the two summands of $(a_{[m-1]})=(a_mu)\cup x_m$ gives
$y_m=(\bar u)\cup z_m$ and $(a_{[m-1]})=(\bar u)\cup(\bar u)\cup z_m=0$ (as $-1\in
k^{\times2}$), a contradiction. If $\varepsilon=0$, for $i<m$ the second summands give
$(a_{[m-1]\smallsetminus i})=(\bar u)\cup z_i$ for all $i\in[m-1]$, contradicting the
induction hypothesis with $g=\bar u$. So no quadratic extension lowers $\ed$ by two;
with (a), (b) this gives $\edd2=\edd3=n-1$, $d_{n-1}=2$, $d_{n-2}=4$.
\end{proof}

\begin{remark}
Whether $\edd{2^j}(\tau_\gen)=n-j$ for $2\le j\le n-2$ is open; the first case is
$n=4$: is there a quartic extension of $k(a_1,\dots,a_4)$ splitting all six quaternion
algebras $(a_i,a_j)$? Non-Galois quartics have no kernel description as clean as
Arason's, which is what makes $\edd4$ hard. [open]
\end{remark}

\begin{example}[$\Z/p$ over $\Q$]\label{ex:Zp}
Let $p$ be an odd prime and $\tau$ any $\Z/p$-torsor over $F\supseteq\Q$. By
Corollary~\ref{cor:rou} and Example~\ref{ex:cyc}, $\edd{[F(\mu_p):F]}_\Q(\tau)\le1$,
so $d_1(\tau)\le p-1$ for every such torsor, although $\ed_\Q(\Z/p)\ge2$ for
$p\ge5$ \cite{Le07}: $\ed_k(G)=1$ in characteristic $0$ iff $G\hookrightarrow
\mathrm{PGL}_2(k)$, and $\Z/p\not\hookrightarrow\mathrm{PGL}_2(\Q)$ for $p\ge5$
since $\zeta_p+\zeta_p^{-1}\notin\Q$. The profile of the generic $\Z/p$-torsor over
$\Q$ thus drops to $1$ at a degree dividing $p-1$, prime to $p$, consistent with
Theorem~\ref{thm:div} since $\ed_\Q(\Z/p;p)=1$. Whether $d_1(\tau_\gen)=
[\Q(\zeta_p+\zeta_p^{-1}):\Q]=(p-1)/2$ is open.
\stag{proved from \cite{Le07}}
\end{example}

\begin{example}[symmetric groups]\label{ex:Sn}
Let $\operatorname{char}k=0$, $L=k(x_1,\dots,x_n)$, $F=K=k(s_1,\dots,s_n)$, so
$\tau_\gen$ is the generic degree-$n$ \'etale algebra. For $F_m=F(x_1,\dots,x_m)$,
$[F_m:F]=n(n-1)\cdots(n-m+1)$, Corollary~\ref{cor:sub} and (B4) give
\[
\ed_k(\tau_\gen\otimes F_m)=\ed_k(S_{n-m}),\qquad\text{so}\qquad
\edd{n(n-1)\cdots(n-m+1)}_k(\tau_\gen)\le\ed_k(S_{n-m}).
\]
With $\ed_k(S_n;q)=\lfloor n/q\rfloor$ \cite{MR}, Theorem~\ref{thm:div} gives
$\edd d(\tau_\gen)\ge\lfloor n/q_d\rfloor$ for the least prime $q_d>d$, so the profile
of the generic \'etale algebra decays no faster than $\sim n/(2d)$, whereas the
root-field bound needs degree $\sim n^m$ for a drop of $m$. For $S_6$ (with
$\ed_k(S_6)=3$, $\ed_k(S_5)=\ed_k(S_4)=2$, $\ed_k(S_3)=1$ \cite{BR}, and
$\ed_k(A_6)=3$ [cited from memory]): $\edd1=3$; $\edd d\ge2$ for $2\le d\le5$;
$\edd6\le2$, $\edd{120}\le1$, $d_0=720$; $2\mid d_2$, $6\mid d_1$; so
$d_2\in\{2,4,6\}$ and $d_1\in6\Z\cap[6,120]$. The discriminant field $L^{A_6}$ is
\emph{not} a witness for $d_2=2$, since over it $\ed=\ed_k(A_6)=3$.
\stag{proved from cited inputs}
\end{example}

\subsection{Provenance: rewrites versus new statements}\label{ss:prov}

\begin{center}
\small
\begin{tabular}{@{}p{0.34\textwidth}p{0.6\textwidth}@{}}
\toprule
Statement & Status relative to the literature\\
\midrule
Prop.~\ref{prop:shape} (shape) & Rewrite of Prop.~3.2 for arbitrary $G$, $k$; formal.\\
Prop.~\ref{prop:d0} (last jump) & New but elementary; (a) is folklore, (b) the
regularity refinement is new.\\
Prop.~\ref{prop:bfe}, Cor.~\ref{cor:rou} & New (elementary); the values quoted are
\cite{BR,KM}.\\
Lemma~\ref{lem:qcl} & Rewrite of \cite[(2.5)]{RS}; the embedding argument is standard.\\
Lemma~\ref{lem:pi} & Rewrite of Prop.~3.6 in a form valid for every finite $G$
(no Witt vectors).\\
Prop.~\ref{prop:genq} & Known: \cite{RS,Me} ($\ed_q$ of a versal object).\\
Thm.~\ref{thm:div} & The mechanism is \cite[Thm.~10.1]{RS}; the statement about jump
degrees and the bound (c) are new.\\
Prop.~\ref{prop:qgp} & New corollary of \cite{KM}.\\
Thm.~\ref{thm:geom} & Rewrite of Thm.~4.2 in full generality; ingredients
\cite{DR,RS}.\\
Lemma~\ref{lem:ind}, Cor.~\ref{cor:sub} & Extends Prop.~3.9 to non-abelian $G$;
probably known to experts, we found no reference.\\
Ex.~\ref{ex:cyc} & Folklore.\\
Thm.~\ref{thm:ea} & (a),(b),(d) are immediate from the above; (c), i.e.\
$\edd2=\edd3=n-1$ and $d_{n-2}=4$, is new (Arason sequence \cite{Ar} + residue
induction).\\
Ex.~\ref{ex:Zp} & New observation from \cite{Le07}.\\
Ex.~\ref{ex:Sn} & The identity $\ed(\tau_\gen\otimes F_m)=\ed(S_{n-m})$ is new but
easy; inputs \cite{BR,MR}.\\
\bottomrule
\end{tabular}
\end{center}

\subsection{Discussion and open questions}
In characteristic $0$ the profile of $\tau_\gen$ is squeezed between Sylow
divisibility (Theorem~\ref{thm:div}: a prime $q$ with $\ed_k(G;q)>j$ divides every
$d_i$, $i\le j$, and the prime-to-$q$ tail is constant), the roots-of-unity layer
(Corollary~\ref{cor:rou}: within degree $[k(\mu_e):k]$, prime to $|G|$ when $|G|$ is
a prime power, the profile drops to the representation-theoretic value
$\ed_{k(\mu_e)}(G)$), and the subgroup upper bounds (Corollary~\ref{cor:sub}). What is
unknown is whether non-Galois or non-subgroup extensions can beat the subgroup
witnesses. In characteristic $p$ for a $p$-group (Sections 3--9) there is no Sylow
divisibility beyond $p\mid d_0$ and no roots-of-unity layer; the whole profile lives in
$p$-power degrees and records Artin--Schreier--Witt layers. Theorem~\ref{thm:geom}
places both regimes in one framework: the profile is the minimal degree of a
$G$-correspondence into a free $G$-variety of bounded dimension, Theorem~\ref{thm:div}
isolates the part of that degree forced by \cite[Thm.~10.1]{RS}, and the rest, the
$q$-power part in characteristic $\ne q$ and the $p$-power part in characteristic $p$,
is where the profile carries information beyond $\ed$ and $\ed(\cdot;q)$.

\begin{question}
(1) For $(\Z/2)^n$ over an algebraically closed field of characteristic $\ne2$, is
$d_j(\tau_\gen)=2^{n-j}$ for all $j$ (known for $j\in\{0,1,n-1,n\}$)? For $(\Z/q)^n$
with $\mu_q\subset k$, the cyclic-kernel analogue of the Arason sequence should give
$d_{n-1}=q$ [expected].
(2) For $\Z/p$ over $\Q$, is $d_1(\tau_\gen)=(p-1)/2$? More generally, when is
Corollary~\ref{cor:rou} sharp, and can an extension not of the form $Fk'$ do better
than adjoining roots of unity?
(3) Determine $d_2(\tau_\gen(S_6))\in\{2,4,6\}$; does $d_j(\tau_\gen(S_n))$ grow
linearly or super-polynomially in $n$?
(4) Is the profile of $\tau_\gen(V)$ independent of the faithful representation $V$?
(Presumably yes, by specialisation as in \cite{RS}.)
\end{question}

% ---------------------------------------------------------------------------
% Section 11: supplements to Sections 6, 7 and 9 (p = 2, n = 2 and n = 3).
% Placement in the source of the paper:
%   11.1 -> after Theorem 6.2 (it settles the "[open]" item there and the p = 2
%           part of Question 6.4);
%   11.2 -> after Proposition 3.7 / Remark 4.5 (Riemann-Roch bound), and its
%           corollary at (2,3) after Theorem 7.1;
%   11.3 -> after Remark 7.6 (Case A empty; U_3 never a quadratic witness; new
%           families with ed = 2; structure of a witness);
%   11.4 -> after Proposition 7.5 / in Section 9 (the cubic slot);
%   11.5 -> replaces the status table of Theorem 7.1 and updates Section 9.2.
% Appendices A-E hold the supporting material.
% ---------------------------------------------------------------------------
\section{Supplements at \texorpdfstring{$(p,n)=(2,2)$ and $(2,3)$}{(p,n)=(2,2) and (2,3)}:
the cubic slot, quadratic witnesses, and a Riemann--Roch bound}\label{sec:supp}

\noindent\textit{Addendum (27 September 2026).} This section settles the open item
of Theorem~6.2, gives a general Riemann--Roch bound for the prime-to-$p$ jump degree,
shows that the Artin--Schreier--Witt curve never carries a quadratic witness at $(2,3)$
and that a quadratic witness can never have a descent field inside $K$, produces new
families of quadratic extensions with essential dimension exactly $2$, and reformulates
the cubic slot. Notation is that of Sections 2, 5--7: $k$ algebraically closed of
characteristic $2$, $K=k(s_0,s_1,s_2)$, $L=k(z_0,z_1,z_2)$, $s=\wp(z)$, $K_2=k(s_0,s_1)$,
$\tau_2$ the generic $\Z/4$-torsor, $E'\colon w^2+w=u(u+1)(u+s_0)+s_1$, $\bar U_n$ the
Artin--Schreier--Witt curve $\wp y=[t]$, $\Phi_2$ the polynomial of Proposition~7.5 (written
out in \eqref{eq:Phi2}). We write $(s_0,s_1)_K$ for the $\Z/4$-class $(s_0,s_1)\in W_2(K)/\wp$,
i.e.\ $\tau_2\otimes_{K_2}K$; by Lemma~4.1(i), $\ed_k((s_0,s_1)_K)=\ed_k(\tau_2)=2$. Every
separable quadratic extension of $K$ is $K(v_f)$, $v_f^2+v_f=f$, $f\in K/\wp K$
(Proposition~3.6). Scripts referred to are listed in Appendix~\ref{app:scripts}.

\subsection{Supplement to Theorem 6.2: the prime-to-\texorpdfstring{$2$}{2} jump degree at
\texorpdfstring{$n=2$}{n=2}}

\begin{theorem}\label{thm:A}
Let $p=2$, $n=2$. Let $u$ be a root of
\[
P(u):=u^3+(1+s_0)u^2+s_0u+s_1\;=\;u(u+1)(u+s_0)+s_1 .
\]
Then $F'=K_2(u)$ is a separable cubic extension of $K_2$ and
$(s_0,s_1)\equiv[s_0+u^2+u]$ mod $\wp W_2(F')$. Hence $\ed_k(\tau_2\otimes F')=1$ and
\[
d_1^{(p')}(\tau_2)=3 .
\]
More precisely, the cubic closed points of $E'$ over $K_2$ are exactly the sections by
the lines $w=\alpha u+\beta$, $\alpha,\beta\in K_2$; each is a single free point of degree
$3$, and over its residue field $(s_0,s_1)$ is Teichm\"uller. Consequently
$D(U_2)=\Z_{\ge2}$, and the profile of $\tau_2$ is completely known:
$\edd1=2$, $\edd2=\edd3=1$, $\edd4=0$, $d_1=2$, $d_1^{(p')}=3$, $d_0=4$.
\stag{proved}
\end{theorem}

\begin{proof}
$P\in k(s_0)[u][s_1]$ is of degree $1$ in $s_1$ with coprime coefficients $1$ and
$u(u+1)(u+s_0)$, hence irreducible in $k(s_0)[u,s_1]$, hence (Gauss, $P$ monic in $u$)
irreducible in $K_2[u]$; an irreducible cubic in characteristic $2$ is separable. By
(3)--(4), for $h=(u,0)\in W_2(F')$: $(s+\wp h)_0=s_0+u^2+u$ and
$(s+\wp h)_1=s_1+u^3+u^2+s_0u^2+s_0u=P(u)=0$, so $s+\wp h=[s_0+u^2+u]$; equivalently
$(u,0)$ is an $F'$-point of $E'^\circ=\tau_2U_2$ and the last sentence of Theorem~6.2
applies. So $\ed_k(\tau_2\otimes F')\le1$, and $\ge1$ since $3<4=d_0$. As
$d_1^{(p')}$ is odd and $\ge d_1(\tau_2)=2$ (Proposition~6.1), $d_1^{(p')}(\tau_2)=3$.

Cubic points: $E'(K_2)=\{\infty'\}$ (Theorem~6.2), so $\operatorname{Pic}^0(E')(K_2)=0$
and every effective $K_2$-divisor of degree $3$ is linearly equivalent to $3\infty'$,
i.e.\ is the zero divisor of an element of $L(3\infty')=\langle1,u,w\rangle$; the
sections $u=c$ are a degree-$2$ fibre plus $\infty'$, so the cubic points are the
sections by lines $w=\alpha u+\beta$. Such a section is a single closed point of degree
$3$ because a $K_2$-rational component would be a $K_2$-point of $E'^\circ$. It lies in
$E'^\circ$, hence is free. $D(U_2)\supseteq\{2,3\}$ and Appendix~\ref{app:genus} gives
all of $\Z_{\ge2}$.
\end{proof}

\begin{remark}
(i) This closes the item marked [open] in Theorem~6.2 and the $p=2$ part of
Question~6.4: the ``group-theoretic'' obstruction $E'(K_2)=0$ kills degree-$1$ points and
\emph{produces} the cubic ones. (ii) For odd $p$ the same mechanism gives
$d_1^{(p')}(\tau_2)\le p^2-p+1$ (Theorem~\ref{thm:B}); Question~6.4 for degrees in
$[2,p-1]$ is untouched.
\end{remark}

\subsection{Supplement to Proposition 3.7 and Remark 4.5: a Riemann--Roch bound}

Let $l_n=l_n(p)$ be the largest lower ramification break of $\bar U_n\to\mathbb P^1_t$ at
$\infty$, i.e.\ the pole order at $\infty$ of the reduced top Witt coordinate of $U_n$;
by Appendix~\ref{app:genus},
\begin{equation}\label{eq:ln}
l_n(p)=1+(p-1)\sum_{j=1}^{n-1}p^{2j-1}:\qquad
l_2(2)=3,\quad l_2(p)=p^2-p+1,\quad l_3(2)=11,\quad l_4(2)=43 .
\end{equation}

\begin{theorem}\label{thm:B}
For every $(p,n)$, $D(U_n)$ contains an integer prime to $p$ and $\le l_n$; hence
\[
d_1^{(p')}(\tau_\gen)\le l_n(p).
\]
This is the best bound obtainable from the pencils $|d\cdot\infty''|$ on
$\tau_\gen\bar U_n$: every function regular on the affine curve $U_n$ with pole order prime
to $p$ has pole order $\ge l_n$.
\stag{proved; the closed formula \eqref{eq:ln} is proved for $(2,2)$, $(2,3)$, $(p,2)$
and cited in general}
\end{theorem}

\begin{proof}
$X:=\tau_\gen\bar U_n$ is a smooth projective geometrically integral $K$-curve with the
$K$-point $\infty''$ (twist of the $G$-fixed point), and its free locus is
$X\smallsetminus\{\infty''\}$ (Proposition~3.7(c)). Riemann--Roch dimensions
$\ell(d\infty'')$ are invariant under base change to $\bar K$, where
$(X,\infty'')\cong(\bar U_n,\infty)$; so the Weierstrass semigroup of $\infty''$ on $X$ is
that of $\infty$ on $\bar U_n$. The top layer $U_n\to U_{n-1}$ is $y^p-y=c$ with
$v_\infty(c)=-l_n$ on $U_{n-1}$ and $\infty$ totally ramified, so $v(y)=-l_n$ on $U_n$
and $l_n$ is a pole order. Hence there is $f\in K(X)$ with pole divisor exactly
$l_n\infty''$. For $c\in K$ the divisor $\operatorname{div}(f-c)+l_n\infty''$ is
effective, $K$-rational, of degree $l_n$ and supported on $X\smallsetminus\{\infty''\}$.
Its closed points have degrees summing to $l_n\equiv1\pmod p$, so one of them has degree
$e$ prime to $p$, $e\le l_n$, and it is free; Theorem~4.2 gives $\edd e_k(\tau_\gen)\le1$.
For optimality: the coordinate ring of $U_n$ is free over that of $U_{n-1}$ with basis
$1,y,\dots,y^{p-1}$, and $v(a_jy^j)=p\,v_{U_{n-1}}(a_j)-jl_n$ are pairwise distinct
modulo $p$, so a function with pole order prime to $p$ has a dominant term with $j\ge1$
and pole order $\ge l_n$.
\end{proof}

\begin{corollary}\label{cor:B23}
At $(p,n)=(2,3)$: the fibres of the pole-order-$11$ pencil are free divisors of degree
$11$ without degree-$1$ components (Proposition~7.5), so each contains a free point of
odd degree in $\{3,5,7,9,11\}$, and $11\in D(U_3)$ exactly (Appendix~\ref{app:genus}).
Hence
\[
d_1^{(p')}(\tau_\gen)\in\{3,5,7,9,11\},\qquad \delta_{\rm free}(U_3)\in\{3,5,7,9,11\},
\]
the second by Corollary~\ref{cor:noquad} below. This refines ``$d_1^{(p')}\ge3$ finite''
in Theorem~7.1. \stag{proved}
\end{corollary}

\begin{remark}
The naive argument ``$\ell(d\infty'')\ge2$ for $d\ge g+1$, take $d$ odd'' does not
work: the fibres of $f\in L(d\infty'')$ have degree equal to the actual pole order of
$f$, which lies in the Weierstrass semigroup and may be even. At $(2,3)$, $g(\bar U_3)=7$
and the semigroup of $\infty$ is $\langle4,6,11\rangle$, whose gaps are
$\{1,2,3,5,7,9,13\}$: every function of pole order $\le10$ has even pole order.
\end{remark}

\subsection{Supplement to Section 7: quadratic witnesses}

Fix $f\in K\smallsetminus\wp K$ and $K'=K(v_f)$. A \emph{witness} is such a $K'$ with
$\ed_k(\tau_\gen\otimes K')=1$; a \emph{descent field} is a subfield $k\subset F_0\subset
K'$ of transcendence degree $1$ with a class $\beta\in W_3(F_0)/\wp$, $s\equiv\beta$ in
$W_3(K')/\wp$.

\begin{theorem}[no descent field inside $K$]\label{thm:caseA}
Let $K'=K(v_f)$ be a witness with descent field $F_0$ and class $\beta$.
\begin{enumerate}
\item[(a)] If $F_0\subset K$ then $F_0=k(t)$ is rational, $s\equiv\beta+\varepsilon(0,0,f)$
mod $\wp W_3(K)$ with $\varepsilon\in\{0,1\}$, and $\ed_k(\tau_\gen)\le\tr_kk(t,f)\le2$.
\item[(b)] The case (a) does not occur: $F_0\not\subset K$.
\item[(c)] The truncation of $\beta$ is a descent of $(s_0,s_1)_K\otimes K'$, so
$f\in\mathfrak F:=\{f\in K/\wp K:\ed_k((s_0,s_1)_K\otimes K(v_f))=1\}$.
\end{enumerate}
\stag{proved}
\end{theorem}

\begin{proof}
(a) $\ker(H^1(K,\Z/8)\to H^1(K',\Z/8))=\operatorname{Hom}(\Z/2,\Z/8)\cong\Z/2$ by
inflation--restriction, generated by the induced class $V^2[f]=(0,0,f)$ (Section~2.2).
Since $s$ and $\beta$ both lie in $W_3(K)/\wp$ and agree over $K'$, $s\equiv\beta+
\varepsilon(0,0,f)=(\beta_0,\beta_1,\beta_2+\varepsilon f)$ by (2); and a subfield of
transcendence degree $1$ of $k(s_0,s_1,s_2)$ is $k(t)$ (a curve dominated by $\A^3$ is
dominated by a line, Lemma~2.5(i)). (b) Truncate the congruence of (a) to $W_2$: the
term $(0,0,f)$ dies and $(s_0,s_1)\equiv(\beta_0,\beta_1)$ mod $\wp W_2(K)$ with
$\beta_0,\beta_1\in F_0$, so $\ed_k((s_0,s_1)_K)\le1$, contradicting
$\ed_k((s_0,s_1)_K)=2$. (c) Truncation commutes with $\wp$ and with base change.
\end{proof}

So a quadratic witness, if one exists, has a descent field \emph{not} contained in $K$
and can never be turned into a counterexample to Ledet's conjecture by this route; the
truncation argument is the one used in the proof of Proposition~7.5 (the step
``$\varepsilon=1\Rightarrow$ a $K$-point of $E'$''), applied to an arbitrary descent
class instead of $[t]$.

\begin{corollary}[$\bar U_3$ admits no quadratic witness]\label{cor:noquad}
For every quadratic $K'/K$, $\tau_\gen\bar U_3$ has no free $K'$-point: $s$ is not
Teichm\"uller over any quadratic extension of $K$, and $D(U_3)\cap\{1,2\}=\emptyset$.
In particular Algorithm~9.1 has nothing to find, and the ``points not coming from the
$s_2$-line'' of Remark~7.6 do not exist. \stag{proved}
\end{corollary}

\begin{proof}
By Proposition~7.5 a free $K'$-point needs $Q=(u,w_1)\in E'(K')$ with $u\notin K$. Let
$\sigma$ generate $\operatorname{Gal}(K'/K)$. $Q+\sigma Q$ is $\sigma$-invariant, hence
lies in $E'(K)=\{\infty'\}$ (Lemma~4.1(ii)); so $\sigma Q=-Q=(u,w_1+1)$, whence
$\sigma(u)=u$ and $u\in K$.
\end{proof}

The same argument at level $2$ shows that $E'(K(v_f))\ne\{\infty'\}$ iff $f\equiv
s_1+u(u+1)(u+s_0)$ mod $\wp K$ for some $u\in K$: the $E_0$-witnesses of Theorem~6.2
are exactly the $u$-fibres (Proposition~\ref{prop:FE0}).

\begin{proposition}[fixed-field trick]\label{prop:fft}
Let $K'=K(v_f)$ with $f\notin\{0,s_0\}+\wp K$, so that $L'=L\otimes_KK'=L(v_f)$ is a
field. Suppose there are a subfield $M\subset L'$ and $x\in L'$ with $L'=M(x)$ purely
transcendental, such that $M$ is fixed pointwise by some $1\ne h\in G$. Then
$\ed_k(\tau_\gen\otimes K')\ge2$. The same holds for $(s_0,s_1)_K\otimes K'$, with $L'$
replaced by $k(z_0,z_1,s_2)(v_f)$ and $G$ by $\Z/4$. \stag{proved}
\end{proposition}

\begin{proof}
A witness gives, by Corollary~4.4, a smooth projective curve $C$ with faithful
$G$-action and a $G$-equivariant dominant rational map from a model of $L'$; take the
model $B\times\A^1$ with $k(B)=M$. Lemma~7.2 for $B\times\A^1\to B$ gives: either a fibre
$\A^1$ maps nonconstantly to $C$, so $C\cong\mathbb P^1$ (Lemma~2.5(i)), impossible
(Lemma~2.5(ii)); or $k(C)\subset M\subset(L')^h$, so $h$ acts trivially on $C$,
contradicting faithfulness.
\end{proof}

\begin{theorem}[quadratic extensions that are not witnesses]\label{thm:fam}
$\ed_k(\tau_\gen\otimes K(v_f))\ge2$, i.e.\ $K(v_f)$ is not a witness, for every
$f\notin\wp K$ in each of the classes
\begin{enumerate}
\item[(a)] $f\in K_2+\wp K$ (this contains $K(v_1)$ and $K(z_0)$ of Theorem~7.3);
\item[(b)] $f\in s_2+K_2+\wp K$ (this contains $K(v_2)$, not covered by Theorem~7.3);
\item[(c)] $f\in k(g)+\wp K$ for any $g\in K$ with $K=K_2(g)$, e.g.\ $f\in k(s_2)$,
$f=1/(s_2+s_0s_1)$;
\item[(d)] $f\in k(s_0,s_1)[s_2]$ a polynomial in $s_2$ of degree $\le7$.
\end{enumerate}
In cases (a), (c) already $\ed_k((s_0,s_1)_K\otimes K(v_f))=2$, i.e.\ $f\notin\mathfrak F$.
Equality $\ed_k(\tau_\gen\otimes K(v_f))=2$ holds in case (b), and in case (a) for
$f\equiv s_0$ or $f\equiv s_1$ (Theorem~7.3); for the other $f$ the value is $2$ or
$\ed_k(\Z/8)$, and only the lower bound is claimed.
\stag{proved}
\end{theorem}

\begin{proof}
In case (b) the upper bound is Theorem~7.1 (the class $(s_0,s_1,e)$ has a descent of
transcendence degree $2$); in general $\ed_k(\tau_\gen\otimes K')\le\ed_k(\tau_\gen)$. (a) Replace $f$ by $e\in K_2$ with $f\equiv e$. If
$e\equiv s_0$, Proposition~3.9. Otherwise $M:=k(z_0,z_1,v_e)$ is algebraic over
$k(z_0,z_1)$, $L'=M(z_2)$, and $g^4\colon z\mapsto z+V^2(1)=(z_0,z_1,z_2+1)$ fixes $M$;
apply Proposition~\ref{prop:fft}. (b) Write $f=s_2+e+\wp h$, $e\in K_2$. Over
$K'=K(v_f)$, $s_2=\wp(v_f+h)+e$, so $s=(s_0,s_1,e)+\wp V^2(v_f+h)\equiv(s_0,s_1,e)$ by
(2). Now $K'=K_2(v_f)$ is purely transcendental over $K_2$, so by Lemma~4.1(i)
$\ed_k(\tau_\gen\otimes K')=\ed_k((s_0,s_1,e)_{K_2})\ge\ed_k((s_0,s_1)_{K_2})=2$.
(c) Replace $f$ by an element of $k(g)$. Put $F:=k(g,v_f)$; then $K'=F(s_0,s_1)$ with
$s_0,s_1$ algebraically independent over $F$, and the total space of $(s_0,s_1)_K\otimes
K'$ is $k(z_0,z_1,g,v_f)=F(z_0,z_1)$, a field on which $\Z/4$ acts fixing $F$; the
argument of Proposition~\ref{prop:fft} with the fibration $B\times\A^2\to B$ excludes a
faithful $\Z/4$-curve, so $\ed_k((s_0,s_1)_K\otimes K')=2$, and a descent of $s$ over
$K'$ would truncate to a descent of $(s_0,s_1)_K\otimes K'$. (d) If $f\in K_2+\wp K$ use
(a); otherwise Theorem~\ref{thm:struct}(b) gives a witness curve of $p$-rank $0$ and
genus $\le6$, contradicting Corollary~\ref{cor:g7}.
\end{proof}

\begin{theorem}[structure of a quadratic witness]\label{thm:struct}
Let $K'=K(v_f)$ be a witness with $f\notin K_2+\wp K$, $C$ its curve, $C'=C/G$. For
$a\in k^2$ let $D_{f,a}\colon y^2+y=f(a_0,a_1,s_2)$ (a curve over $k$ in the $s_2$-line)
and, for $b\in\wp_2^{-1}(a)\subset\A^2_z$, let $D_b$ be its pullback along
$z_2\mapsto s_2=z_2^2+z_2+C(b_0,b_1)$.
\begin{enumerate}
\item[(a)] For $a$ in a dense open subset of $\A^2(k)$: $D_{f,a}$ is integral and
dominates $C'$, and every irreducible component of $D_b$ dominates $C$. Hence
$g(C)\le g(\hat D_b)$ and $\sigma(C)\le\sigma(\hat D_b)$ ($\sigma$ the $p$-rank, $\hat D_b$
the normalisation of a component).
\item[(b)] If $f\in k(s_0,s_1)[s_2]$ has $s_2$-degree $m\ge1$, then $\sigma(C)=0$,
$C\to C'$ is totally ramified at exactly one point and \'etale elsewhere, $C'$ has
$p$-rank $0$ and genus $\le\lfloor(m-1)/2\rfloor$, and $7\le g(C)\le m-1$. In particular
$m\ge8$, and for $8\le m\le15$ one has $C'\cong\mathbb P^1$: $C$ is the
Artin--Schreier--Witt curve of a polynomial Witt vector $\beta\in W_3(k[t])$, and
$s\equiv\beta(t)$ over $K'$ with $t\in K'\smallsetminus K$.
\end{enumerate}
\stag{proved}
\end{theorem}

\begin{proof}
(a) Let $B_f\colon y^2+y=f(s_0,s_1,s_2)$ over $\A^3_s$, $k(B_f)=K'$. Its generic fibre
over $\A^2_{(s_0,s_1)}$ is the Artin--Schreier curve $y^2+y=f$ over $K_2$ in the
$s_2$-line, geometrically integral because $f\notin K_2+\wp K$ (a constant quadratic
subextension $K_2(v_e)\subset K'$ would force $f\equiv e$). Apply Lemma~7.2 to $B_f\to
\A^2_s$ and the dominant rational map $B_f\dashrightarrow C'$ given by $k(C')\subset K'$:
the alternative $k(C')\subset k(s_0,s_1)\subset K$ is excluded by
Theorem~\ref{thm:caseA}(b), so $D_{f,a}\dashrightarrow C'$ is dominant for general $a$.
Let $Y_f$ be the normalisation of $\A^3_z\times_{\A^3_s}B_f$, with function field $L'$;
the equivariant map $Y_f\dashrightarrow C$ composed with $C\to C'$ factors through
$Y_f\to B_f\dashrightarrow C'$, and the fibre of $Y_f$ over $a$ is
$\bigsqcup_{b\in\wp_2^{-1}(a)}D_b$, permuted transitively by $G$. Each $D_b\to D_{f,a}
\to C'$ is nonconstant, hence so is $D_b\to C$. Genus and $p$-rank do not increase under
dominant maps of curves. (b) $D_{f,a}$ is $y^2+y=$ (polynomial of degree $m$ in $s_2$),
with a single branch point, so $\sigma(D_{f,a})=0$ (Deuring--Shafarevich) and
$g\le\lfloor(m-1)/2\rfloor$; $D_b$ is $y^2+y=$ (polynomial of degree $2m$ in $z_2$),
again with one branch point, so $\sigma=0$ and, after Artin--Schreier reduction to odd
degree $\le2m-1$, $g(D_b)\le m-1$. By (a), $\sigma(C)=\sigma(C')=0$ and $g(C)\le m-1$;
Lemma~\ref{lem:prank0} and Corollary~\ref{cor:g7} give the ramification structure and
$g(C)\ge7$. If $g(C')=1$ then Riemann--Hurwitz gives $2g(C)-2\ge0+28$, so $m\ge16$; for
$m\le15$, $C'=\mathbb P^1$ and the cover is \'etale outside one point, so its class
lies in $H^1(\A^1,\Z/8)=W_3(k[t])/\wp$, with $\beta_0$ nonconstant since
$s_0\notin\wp K'$ (else $K'=K(z_0)$).
\end{proof}

\begin{corollary}[minimal genus]\label{cor:g7}
A smooth projective curve with faithful $\Z/8$-action and $p$-rank $0$ in
characteristic $2$ has genus $\ge7$, with equality iff $C/G\cong\mathbb P^1$ and the
upper ramification breaks at the unique ramified point are $(1,2,4)$; $\bar U_3$ is an
example, $g(\bar U_3)=7$. A faithful $\Z/4$-curve of $p$-rank $0$ has genus $\ge1$, with
equality iff it is $(E_0,\text{Witt action})$ of Theorem~6.2. \stag{proved}
\end{corollary}

\begin{proof}
Lemma~\ref{lem:prank0}: one totally ramified point, $C/G$ of $p$-rank $0$. The different
exponent at that point is $(u_1+1)+2(u_2+1)+4(u_3+1)$ in terms of the upper breaks, with
$u_1\ge1$ and $u_{i+1}\ge2u_i$ \cite{Ga,Th}, so $2g-2\ge-16+28=12$. For $\Z/4$:
$2g-2\ge-8+2+6=0$, and the genus-$1$ case is $(E_0,\text{Witt})$ by
Lemma~\ref{lem:Z4ss}.
\end{proof}

\begin{remark}
(i) Theorem~7.3's proofs for $K(v_1)$ and $K(v_1,v_2)$ are instances of the mechanism
of Theorem~\ref{thm:struct}; for $f\in K_2$ the fixed-field trick is shorter, but for
the quartic tower it does not apply directly ($g^4$ moves $y'=v_2+z_2$) and the argument
of Theorem~7.3 is needed. (ii) The only genus-$1$ faithful $\Z/8$-curves are ordinary
elliptic curves with $G$ acting by translation by a point of order $8$
(Appendix~\ref{app:RH}); by Theorem~\ref{thm:struct}(a) they can occur as witness curves
only if $f$ has, for general $a$, at least two poles on the $s_2$-line, and then
$\tau_\gen\otimes K'$ is the \'etale $\Z/8$-isogeny torsor of an ordinary elliptic curve
at a non-constant point (Appendix~\ref{app:isog}). This and polynomial $f$ of
$s_2$-degree $\ge8$ are the only remaining shapes of a quadratic witness.
\end{remark}

\subsection{Supplement to Proposition 7.5 and Question (Q2): the cubic slot}

Let $F'/K$ be cubic (separable without loss, Proposition~3.6). Then $L\otimes_KF'$ is a
field and Corollary~4.4 applies: a witness for $\edd3(\tau_\gen)=1$ is a faithful
$\Z/8$-curve $C$ with a free $F'$-point on $\tau_\gen C$.

\begin{proposition}\label{prop:cubiccurves}
A cubic witness curve $C$ has genus $\ge3$ and, if its $p$-rank is $0$, genus $\ge7$.
Ordinary elliptic curves with translation action carry no free point of odd degree.
\stag{proved}
\end{proposition}

\begin{proof}
$\mathbb P^1$ and supersingular elliptic curves are excluded by Lemma~2.5. If $C$ is an
ordinary elliptic curve with $G$ acting by translation by $P$ of order $8$ (the only
genus-$1$ case, Appendix~\ref{app:RH}), then $\tau_\gen C$ is a principal homogeneous
space under $C$ whose period divides $8$; period and index have the same prime divisors
\cite{LT}, so a nontrivial such torsor has no closed point of odd degree, and a trivial
one has a $K$-point, free since the action is free, forcing $\ed_k(\tau_\gen)\le1$,
false. Genus $2$: the hyperelliptic involution is central, and the image of $\Z/8$ in
$\operatorname{Aut}(\mathbb P^1)$ is $\Z/8$ or $\Z/4$, excluded by Lemma~2.5(ii).
$p$-rank $0$: Corollary~\ref{cor:g7}.
\end{proof}

On $\bar U_3$ the question is completely explicit. By Proposition~5.2(a) and (3)--(4), a
free $F'$-point of $\tau_\gen\bar U_3$ is $h=(u,w_1,w_2)\in W_3(F')$ with
\begin{equation}\label{eq:U3eq}
w_1^2+w_1=u(u+1)(u+s_0)+s_1,\qquad w_2^2+w_2=e_Q:=s_2+\Phi_2(s_0,s_1,u,w_1),
\end{equation}
and then $s\equiv[s_0+u^2+u]$ over $F'$.

\begin{proposition}[reformulation]\label{prop:C}
$\edd3(\tau_\gen)=1$ is witnessed on $\bar U_3$ iff there are $\alpha,\beta\in K$ such
that, for $u$ a root of
\[
u^3+(1+s_0+\alpha^2)u^2+(s_0+\alpha)u+(s_1+\beta^2+\beta)=0,
\]
one has $e_Q=s_2+\Phi_2(s_0,s_1,u,\alpha u+\beta)\in\wp K(u)$. Every such cubic is
irreducible over $K$, and these are all the cubic points of $E'$ over $K$.
\stag{proved}
\end{proposition}

\begin{proof}
As in Theorem~\ref{thm:A}, $E'(K)=\{\infty'\}$ (Lemma~4.1(ii)) gives
$\operatorname{Pic}^0(E')(K)=0$, so the cubic points of $E'$ over $K$ are the line
sections, each irreducible. A free $F'$-point with $[F':K]=3$ has $F'=K(u)$: if $u\in K$
then $t=s_0+u^2+u\in K$ and $s\equiv[t]$ over $F'$ would give $s\equiv[t]$ over $K$,
since $\res_{F'/K}$ is injective on $H^1(\cdot,\Z/8)$ ($\cores\circ\res=3$ is a unit
mod $8$).
\end{proof}

\begin{proposition}[necessary condition and a partial obstruction]\label{prop:D}
\begin{enumerate}
\item[(a)] If $s\equiv[t]$ over a cubic $F'=K(u)$ then $s\equiv3\cdot\cores_{F'/K}[t]$ in
$W_3(K)/\wp$, where $\cores$ is the Witt-vector trace
$[t]+[t']+[t'']=(e_1,e_2,e_1e_3+e_1^2e_2)$ ($e_i$ the elementary symmetric functions of
the conjugates). Coordinates $0$ and $1$ of this condition hold automatically for every
line, and coordinate $2$ reads
\begin{equation}\label{eq:N1}
\operatorname{Tr}_{K(u)/K}(e_Q)=s_2+R(s_0,s_1,\alpha,\beta)\in\wp K
\end{equation}
with $R$ the explicit polynomial of Appendix~\ref{app:cubic}.
\item[(b)] \eqref{eq:N1} fails for all $\alpha,\beta\in K_2$, and more generally for all
$\alpha,\beta$ with pole order $\le1$ at $s_2=\infty$. Hence a cubic witness on
$\bar U_3$, if any, has a line with $\alpha$ or $\beta$ of pole order $\ge2$ at
$s_2=\infty$.
\end{enumerate}
\stag{(a) proved, coordinates checked symbolically; (b) proved, with computer-verified
coefficients}
\end{proposition}

\begin{proof}
(a) $\cores\circ\res=3$ on $H^1(K,\Z/8)$ and $3^2\equiv1\pmod8$. Corestriction is a
morphism of $\delta$-functors for $0\to\Z/8\to W_3\xrightarrow{\wp}W_3\to0$, and on
$H^0$ it is the norm, i.e.\ the Witt trace; the formula for $[t]+[t']+[t'']$ follows from
(3), and the coordinate computations are in Appendix~\ref{app:cubic}. (b) For
$\alpha,\beta\in K_2$, $s_2+R$ has a simple pole at $s_2=\infty$ (Remark~7.4). In
general write $\alpha=\sum_{i\le1}a_is_2^i$, $\beta=\sum_{i\le1}b_is_2^i$ in
$K_2((1/s_2))$. Membership in $\wp$ of the completion forces, for every odd $i_0\ge1$, the
chain condition $\sum_{j}c_{i_02^j}^{2^{J-j}}=0$ on the coefficients $c_i$ of $s_2^i$ in
$s_2+R$ (Appendix~\ref{app:cubic}). The only monomial of $R$ producing $s_2^{11}$ from
exponents $\le1$ is $\alpha^{11}$, giving $c_{11}=a_1^{11}$ and $a_1=0$; then
$c_3=b_1^3$ gives $b_1=0$; then $\alpha,\beta$ are integral at $\infty$ and
$s_2+R$ has a simple pole, contradicting the chain condition for $i_0=1$.
\end{proof}

Exhaustive searches over $\mathbb F_2$-polynomial lines (about $2.4\cdot10^6$ pairs
$(\alpha,\beta)$ of small degree, Appendix~\ref{app:cubic}) found no line satisfying
\eqref{eq:N1}. Note that \eqref{eq:N1} is only necessary: the kernel of $\cores$ on
$H^1(K(u),\Z/2)$ is large. Our expectation is that no cubic witness exists on
$\bar U_3$; whether another faithful $\Z/8$-curve of genus $\ge3$ carries one is a
separate question, so (Q2) remains open.

\subsection{Updated status at \texorpdfstring{$(2,3)$}{(2,3)} and revised questions}

\begin{theorem}[status at $(2,3)$]\label{thm:status}
\[
\edd1\in\{2,3\},\quad\edd2\in\{1,2\},\quad\edd3\in\{1,2\},\quad
\edd4=\dots=\edd7=1,\quad\edd8=0;
\]
$d_2\in\{1,2\}$, $d_1\in\{2,3,4\}$, $d_0=8$, $d_1^{(p')}\in\{3,5,7,9,11\}$. Moreover:
a quadratic witness for $d_1=2$ has no descent field inside $K$, is not $\bar U_3$,
is not any $f$ of Theorem~\ref{thm:fam}, and for polynomial $f$ is an
Artin--Schreier--Witt curve of a polynomial Witt vector of $s_2$-degree $\ge8$;
$D(U_3)\cap\{1,2\}=\emptyset$, $D(U_3)\supseteq\langle4,6,11\rangle\smallsetminus\{0\}$,
and $D(U_3)\cap\{3,5,7,9,13\}$ is unknown. \stag{proved}
\end{theorem}

The open questions of Section~9.2 are refined as follows. (Q1): is there
$f\in K/\wp K$, with at least two poles on the $s_2$-line for general $(s_0,s_1)$ or a
polynomial of $s_2$-degree $\ge8$, such that $K(v_f)$ is a witness? The first shape
requires the \'etale $\Z/8$-isogeny torsor of an ordinary elliptic curve
(Appendix~\ref{app:isog}); the level-$2$ version (is $(s_0,s_1)_K\otimes K(v_f)$ the
\'etale $\Z/4$-isogeny torsor of an ordinary elliptic curve at a non-constant point?) is
the natural first step, since a negative answer there kills all such level-$3$
witnesses by truncation. (Q2): does \eqref{eq:N1} fail for all $\alpha,\beta\in K$
(a local analysis at $s_2=\infty$ for pole orders $\ge2$ is the missing piece), and
which of $3,5,7,9$ lie in $D(U_3)$? The splitting type over $K$ of the degree-$11$
fibres of the pole-order-$11$ pencil is a finite computation that would decide the
latter. (Q3), (Q4) are unchanged, except that $d_1^{(p')}(\tau_2)\le p^2-p+1$ for all
$p$ and $=3$ for $p=2$.

\subsection{A quantitative Reichstein--Vistoli theorem}\label{ss:qRV}

Reichstein--Vistoli \cite{RV} prove $\ed_{k,p}(G)\le1$ for every finite $p$-group $G$ in
characteristic $p$, i.e.\ every $G$-torsor becomes one-dimensional over \emph{some}
extension of degree prime to $p$: in the language of the profile, $d_1^{(p')}(\tau)<\infty$
for every $\tau$. The proof passes through the $p$-closure, a large field, and gives
neither a bound on the degree nor uniformity in $\tau$. The Riemann--Roch argument of
Theorem~\ref{thm:B} gives both. Throughout this subsection $k$ is a field of
characteristic $p$, algebraically closed unless said otherwise, $G$ a finite group, and
$V$ a faithful representation of $G$ over $k$ with generic torsor $\tau_\gen$ over
$K=k(V)^G$.

\begin{proposition}[the generic torsor dominates every profile]\label{prop:vers}
Let $\tau$ be a $G$-torsor over an infinite field $F\supseteq k$. Then for every $d\ge1$
\[
\edd d_k(\tau)\le\edd d_k(\tau_\gen),\qquad
\edd d_k{}'(\tau)\le\edd d_k{}'(\tau_\gen)
\]
(the second for the prime-to-$p$ profiles), hence $d_j(\tau)\le d_j(\tau_\gen)$ and
$d_j^{(p')}(\tau)\le d_j^{(p')}(\tau_\gen)$ for all $j$. Consequently
$\sup_\tau\edd d_k(\tau)=\edd d_k(\tau_\gen)$, and the profile of $\tau_\gen$ does not
depend on the faithful representation $V$. \stag{proved}
\end{proposition}

\begin{proof}
Let $F'/K$ realise $\edd d_k(\tau_\gen)=m$, $[F':K]=d'\le d$; by Proposition~3.6 (or
Lemma~\ref{lem:pi}) we may take $F'/K$ separable. By Theorem~\ref{thm:geom} there are a
quasi-projective generically free primitive $G$-variety $Z$ of dimension $\le m$ and a
$G$-equivariant morphism $\operatorname{Spec}(L\otimes_KF')\to Z^{\rm free}$. As in the
proof of Theorem~4.2, let $W'$ be the normalisation of $V$ in $L\otimes_KF'$; then
$\pi\colon W'\to V$ is finite, $G$-equivariant, generically \'etale of degree $d'$, and
the equivariant map $W'\dashrightarrow Z^{\rm free}$ is defined on a dense $G$-stable
open $U$. Choose a dense $G$-stable open $V_1\subseteq V$ over which $\pi$ is \'etale and
$\pi^{-1}(V_1)\subseteq U$. Twisting by $\tau$ (functoriality of twists,
\cite[\S3]{DR}) gives a finite \'etale $F$-morphism ${}^\tau\pi^{-1}(V_1)\to{}^\tau V_1$
of degree $d'$ and an $F$-morphism ${}^\tau\pi^{-1}(V_1)\to{}^\tau Z^{\rm free}$. Now
${}^\tau V$ is an $F$-form of the vector space $V_F$, hence isomorphic to affine space by
Hilbert~90, so ${}^\tau V_1(F)$ is non-empty as $F$ is infinite; pick $x\in{}^\tau V_1(F)$.
Its fibre is $\operatorname{Spec}A$ for an \'etale $F$-algebra $A=\prod_iF_i$ of
dimension $d'$, and each $\operatorname{Spec}F_i\to{}^\tau Z^{\rm free}$ gives a closed
point of ${}^\tau Z$ in the free locus of degree $\le[F_i:F]$. The degrees $[F_i:F]$
sum to $d'\le d$, so $\edd d_k(\tau)\le m$ by Theorem~\ref{thm:geom}; if $p\nmid d'$,
some $[F_i:F]$ is prime to $p$. The last two statements: the supremum is attained at
$\tau_\gen$, and two generic torsors (for $V$ and $V'$) are torsors over infinite
fields, so each dominates the other.
\end{proof}

\begin{definition}
For a smooth projective curve $C$ over $k$ with a faithful $G$-action and a $G$-fixed
$k$-point $\infty$, let $\lambda(C,\infty)$ be the smallest element of the Weierstrass
semigroup $H(\infty)$ that is prime to $p$, and let $\lambda_k(G)$ be the minimum of
$\lambda(C,\infty)$ over all such pairs ($\lambda_k(G)=\infty$ if none exists).
\end{definition}

Such pairs exist for every finite $p$-group and, more generally, for every
$G\cong P\rtimes\Z/m$ with $P$ a $p$-group and $p\nmid m$: by Harbater and Katz--Gabber
\cite{Ha,Ka} there is a connected $G$-cover of $\mathbb P^1$ \'etale over
$\mathbb G_m$ (over $\A^1$ if $G$ is a $p$-group) with inertia group $G$ at $\infty$, and
the point above $\infty$ is then $G$-fixed. Conversely a $G$-fixed point forces $G$ to be
an inertia group in characteristic $p$, hence of this form. For $G=\Z/p^n$ the pair
$(\bar U_n,\infty)$ is defined over $\mathbb F_p$.

\begin{theorem}[quantitative Reichstein--Vistoli]\label{thm:qRV}
\begin{enumerate}
\item[(a)] Let $(C,\infty)$ be as in the definition. For every $G$-torsor $\tau$ over
every field $F\supseteq k$ there is a separable extension $F'/F$ with $p\nmid[F':F]\le
\lambda(C,\infty)$ and $\ed_k(\tau\otimes F')\le1$. Hence
\[
d_1^{(p')}(\tau)\le d_1^{(p')}(\tau_\gen)\le\lambda_k(G)\qquad\text{for every }\tau .
\]
\item[(b)] For $G=\Z/p^n$: $\lambda_k(\Z/p^n)\le l_n(p)=1+(p-1)\sum_{j=1}^{n-1}p^{2j-1}$,
realised by $(\bar U_n,\infty)$, and $\lambda(C,\infty)=l_n(p)$ is the minimum of
$\lambda(C,\infty)$ over all connected $\Z/p^n$-covers $C\to\mathbb P^1$ \'etale over
$\A^1$. So every $\Z/p^n$-torsor over every $F\supseteq k$ becomes one-dimensional over
a prime-to-$p$ extension of degree at most $l_n(p)$; the bound is $1$ for $n=1$,
$p^2-p+1$ for $n=2$, $11$ at $(2,3)$, $43$ at $(2,4)$, and it is attained by
$\tau_\gen$ for $p^n=4$ (Theorem~\ref{thm:A}). Together with Proposition~3.4,
$d_1(\tau)\le\min(p^{n-1},l_n(p))$ for every $\Z/p^n$-torsor.
\item[(c)] For an arbitrary finite $G$ with Sylow $p$-subgroup $G_p$:
$d_1^{(p')}(\tau)\le[G:G_p]\cdot\lambda_k(G_p)$ for every $G$-torsor $\tau$ over every
$F\supseteq k$.
\end{enumerate}
\stag{proved; (b) uses the inequality $u_{i+1}\ge pu_i$ for upper breaks of cyclic
$p^n$-extensions of $k((x))$ \cite{Ga,Th} and the existence results \cite{Ha,Ka}, both
cited}
\end{theorem}

\begin{proof}
(a) If $F$ is finite, $\ed_k(\tau)=0$. Let $F$ be infinite, $X={}^\tau C$, $\infty'\in X(F)$
the twist of $\infty$. The non-free locus $N\subset C$ is a finite $G$-stable set (fixed
points of the nontrivial elements of a group acting faithfully on a curve), so ${}^\tau N$
is a finite set of closed points of $X$. The Weierstrass semigroup of $\infty'$ on $X$
equals $H(\infty)$ (Riemann--Roch dimensions are invariant under base change, and
$(X,\infty')_{\bar F}\cong(C,\infty)_{\bar F}$), so there is $f\in F(X)$ with pole divisor
exactly $\lambda\infty'$, $\lambda=\lambda(C,\infty)$. The fibres $f=c$, $c\in F$, are
pairwise disjoint effective $F$-rational divisors of degree $\lambda$ not containing
$\infty'$; all but finitely many avoid ${}^\tau N$. For such a fibre the degrees of its
closed points sum to $\lambda$, which is prime to $p$, so one closed point has degree
$e$ prime to $p$, $e\le\lambda$, and lies in the free locus; its residue field $F'$ is
separable over $F$ (its degree is prime to $p$), and Theorem~\ref{thm:geom} with
$Z=C^{\rm free}$ gives $\ed_k(\tau\otimes F')\le1$. The chain of inequalities is
Proposition~\ref{prop:vers}.

(b) The upper bound is Theorem~\ref{thm:B}. For a connected $\Z/p^n$-cover
$C\to\mathbb P^1$ \'etale over $\A^1$ the argument in the proof of Theorem~\ref{thm:B}
applies verbatim ($C\to C/\langle g^{p^{n-1}}\rangle$ is an Artin--Schreier cover \'etale
over the affine part, so the coordinate ring of $C\smallsetminus\{\infty\}$ is free with
basis $1,y,\dots,y^{p-1}$), and shows $\lambda(C,\infty)=l_n(C,\infty)$, the last lower
ramification break at $\infty$. Writing the upper breaks as $u_1$ and
$u_{j+1}=pu_j+\delta_j$ with $\delta_j\ge0$ \cite{Ga,Th}, Herbrand's formula gives
$l_n=u_1+\sum_{j=1}^{n-1}p^j\bigl((p-1)u_j+\delta_j\bigr)$, which is minimal for
$u_1=1$ and all $\delta_j=0$, i.e.\ for the upper breaks $(1,p,\dots,p^{n-1})$ of
$\bar U_n$ (Appendix~\ref{app:genus}). Sharpness at $p^n=4$ is Theorem~\ref{thm:A}; the
last claim is Proposition~3.4 applied to an arbitrary torsor ($\edd{p^{n-1}}_k(\tau_a)\le
\tr_kk(a_0)\le1$).

(c) $T/G_p\to\operatorname{Spec}F$ is \'etale of degree $[G:G_p]$, so it has a field
factor $F_1$ with $p\nmid[F_1:F]\le[G:G_p]$; over $F_1$ the torsor has a point of
$T/G_p$, hence $\tau\otimes F_1\cong\Ind_{G_p}^G\tau_p$ for a $G_p$-torsor $\tau_p$ over
$F_1$. By (a) there is $F_2/F_1$ with $p\nmid[F_2:F_1]\le\lambda_k(G_p)$ and
$\ed_k(\tau_p\otimes F_2)\le1$, and $\ed_k(\tau\otimes F_2)=\ed_k(\tau_p\otimes F_2)$ by
Lemma~\ref{lem:ind}.
\end{proof}

\begin{remark}
(i) What is new relative to \cite{RV} and to Proposition~3.7(b),(c): the explicit bound,
its uniformity in $\tau$, the identification of the optimal bound from one-point covers
with a local ramification invariant for cyclic groups, and the versality of the profile
(Proposition~\ref{prop:vers}), which also answers the question whether the profile of
$\tau_\gen$ depends on $V$. The mechanism, a free point on a curve with a fixed point, is
that of \cite[Lemma~9.1]{RS} and Proposition~3.7(c); what replaces largeness is
Riemann--Roch on the twisted curve. (ii) A crude bound valid for every $(C,\infty)$ is
$\lambda(C,\infty)\le$ the least integer prime to $p$ that is $\ge2g(C)$, since
$H(\infty)\supseteq[2g(C),\infty)$; at $(2,3)$ this gives $15$ against $11$.
(iii) Whether $\lambda_k(\Z/p^n)=l_n(p)$ over \emph{all} curves with a fixed point, not
only one-point covers, is open, as is the exact value of $d_1^{(p')}(\tau_\gen)$ for
$n=2$ and $p$ odd (between $2$ and $p^2-p+1$) and at $(2,3)$ (in $\{3,5,7,9,11\}$).
(iv) The uniform statement is a statement about $\sup_\tau$: by
Proposition~\ref{prop:vers} it is equivalent to the bound for $\tau_\gen$ alone.
\end{remark}

% ---------------------------------------------------------------------------
% Appendices: supporting material for Section 11 (tangential results).
% ---------------------------------------------------------------------------
\appendix

\section{Witt-vector computations at \texorpdfstring{$p=2$, $n=3$}{p=2, n=3}}\label{app:witt}

All formulas were re-derived from the ghost components $w_0=x_0$, $w_1=x_0^2+2x_1$,
$w_2=x_0^4+2x_1^2+4x_2$ over $\Z$ and reduced modulo $2$ (script \texttt{witt\_check.py});
they agree with (3)--(4). \stag{computed}

\emph{Negation.} $-(a_0,a_1,a_2)=(a_0,\ a_1+a_0^2,\ a_2+a_1^2+a_0^2a_1+a_0^4)$; in
particular $-x\ne x$, but $\{s-\wp x\}=\{s+\wp x\}$ as $x$ ranges over $W_3(K)$.

\emph{The general translate.} For $x=(x_0,x_1,x_2)\in W_3(K)$,
$s+\wp x=(t,r_1,r_2)$ with
\begin{align*}
t&=s_0+x_0^2+x_0,\\
r_1&=s_1+x_0(x_0+1)(x_0+s_0)+\wp(x_1),\\
r_2&=s_2+\Phi_2(s_0,s_1,x_0,x_1)+\wp(x_2),
\end{align*}
where
\begin{equation}\label{eq:Phi2}
\begin{aligned}
\Phi_2={}&s_0^3(x_0^2+x_0)+s_0s_1(x_0^2+x_0)+s_0(x_0^6+x_0^4)+s_0(x_0^2+x_0)(x_1^2+x_1)\\
&+s_1(x_0^3+x_0^2)+s_1(x_1^2+x_1)+x_0^7+x_0^4+(x_0^3+x_0^2)(x_1^2+x_1)+x_1^3+x_1^2 .
\end{aligned}
\end{equation}
Thus the $u$-fibre family of Theorem~6.2 is the $x_1$-independent part of $r_1$:
$(s+\wp(c,x_1))_1-(c(c+1)(c+s_0)+s_1)=\wp(x_1)$, and
$(s-\wp x)_1-(s+\wp x)_1=x_0^4+x_0^2\in\wp K$.

\section{Genus, ramification and Weierstrass semigroup of \texorpdfstring{$\bar U_n$}{U\_n}}
\label{app:genus}

\emph{Ramification at $\infty$.} $\bar U_n\to\mathbb P^1_t$ is a $\Z/p^n$-cover, \'etale
over $\A^1$ and totally ramified at $\infty$, of Witt class $[t]$ with $v_\infty(t)=-1$.
By the conductor formula for Artin--Schreier--Witt extensions \cite{Th} (upper breaks
$u_j=\max_{i<j}p^{\,j-1-i}m_i$ for reduced pole orders $m_i$) the upper breaks are
$u_j=p^{j-1}$, and Herbrand's function converts them to the lower breaks
\[
l_1=1,\qquad l_{j+1}=l_j+p^j(u_{j+1}-u_j)=l_j+(p-1)p^{2j-1},
\]
which gives \eqref{eq:ln}. \stag{cited} Direct verification at $(2,3)$ \stag{computed,
\texttt{genus\_U3.py}}: the tower $U_3\to U_2=E_0\to\mathbb P^1_{y_0}\to\mathbb P^1_t$
has layers $y_0^2+y_0=t$ (break $1$), $y_1^2+y_1=y_0^3+y_0^2$ (break $3$) and
$y_2^2+y_2=C(y_0,y_1)$ with $C$ from (4). On $E_0$, using $y_1^2=y_1+y_0^3+y_0^2$,
$C=y_0^7+y_0^6+(y_0^3+y_0^2)y_1$, of pole order $14$; one reduction step,
$C+\wp(y_0^2y_1)=y_0^4y_1+y_0^3y_1$, has odd pole order $8+3=11$. So the lower breaks
are $(1,3,11)$ and the upper breaks $(1,2,4)$, as predicted.

\emph{Genus.} With $l_0:=-1$, Riemann--Hurwitz gives
$2g(\bar U_n)-2=-2p^n+\sum_{j=1}^n(p^{\,n-j+1}-1)(l_j-l_{j-1})$, whence
\[
g(\bar U_2)=\tfrac12p(p-1)^2,\qquad g(\bar U_3)=7\ (p=2),\qquad g(\bar U_4)=35\ (p=2),
\qquad g(\bar U_3)=78\ (p=3).
\]
Cross-check at $(2,3)$: the different exponent is $(u_1+1)+2(u_2+1)+4(u_3+1)=2+6+20=28$
and $2g-2=-16+28=12$. \stag{computed}

\emph{Weierstrass semigroup at $(2,3)$.} The functions $y_0$, $y_1$, $y_2+y_0^2y_1$ are
regular on $U_3$ with pole orders $4$, $6$, $11$ at $\infty$; $\langle4,6,11\rangle$ has
exactly the $7=g$ gaps $\{1,2,3,5,7,9,13\}$, so it is the full semigroup of $\infty$.
The smallest odd pole order of a function regular on $U_3$ is therefore $11$.

\emph{Degrees realised on $\bar U_n$.} Let $H$ be the semigroup of $\infty$ on
$\bar U_n$. For $d\in H$ choose $f\in K(\tau_\gen\bar U_n)$ with pole divisor $d\infty''$
and $df\ne0$ (possible for every $d\in H$ prime to $p$, and at $(2,3)$ for every
$d\in\langle4,6,11\rangle$, e.g.\ $y_0$, $y_1$, $y_0^2+y_1$, $y_0y_1$ for $4,6,8,10$).
Then $K(\tau_\gen\bar U_n)$ is a separable extension of degree $d$ of $K(f)$, and since
$K=k(s_0,\dots,s_{n-1})$ is Hilbertian \cite{FJ}, for infinitely many $c\in K$ the fibre
$f=c$ is a single closed point of degree $d$ in the free locus. Hence
\[
D(U_n)\supseteq\{d\in H\smallsetminus\{0\}:d\text{ prime to }p\},\qquad
D(U_3)\supseteq\langle4,6,11\rangle\smallsetminus\{0\},\qquad D(U_2)=\Z_{\ge2}\ (p=2).
\]
With Corollary~\ref{cor:noquad}, $D(U_3)\cap\{1,2\}=\emptyset$; the unknown part of
$D(U_3)$ is $\{3,5,7,9,13\}$, and the odd degrees realised on $\bar U_3$ so far are the
odd elements $\ge11$ of $\langle4,6,11\rangle$ ($13$ is a gap). \stag{proved, modulo
Hilbert irreducibility}

\section{Cyclic \texorpdfstring{$2$}{2}-power actions on curves in characteristic
\texorpdfstring{$2$}{2}}\label{app:RH}

Let $C$ be a smooth projective curve over $k$ with a faithful action of $G=\Z/2^n$,
$\sigma(\cdot)$ the $p$-rank. For a $G$-orbit with inertia of order $2^a$ and upper
breaks $u_1<\dots<u_a$ one has $u_1\ge1$ and $u_{i+1}\ge2u_i$ \cite{Ga,Th}; the different
exponent at a point of that orbit is $\sum_{i=1}^a(2^i-2^{i-1})(u_i+1)$ in terms of the
upper breaks of the inertia group. Deuring--Shafarevich \cite{Cr}:
\[
\sigma(C)-1=2^n(\sigma(C/G)-1)+\sum_{\text{orbits}}2^n\Bigl(1-\frac1{e_{\rm orbit}}\Bigr).
\]

\begin{lemma}\label{lem:prank0}
If $\sigma(C)=0$ then $\sigma(C/G)=0$, and $C\to C/G$ is totally ramified at exactly
one point and \'etale elsewhere. Conversely such a cover of a $p$-rank-$0$ curve has
$p$-rank $0$. \stag{proved}
\end{lemma}

\begin{proof}
$2^n\sigma(C/G)+\sum2^n(1-1/e)=2^n-1$; each orbit with $e\ge2$ contributes at least
$2^{n-1}$, and $2^n-1$ is reached only by a single orbit with $e=2^n$.
\end{proof}

\begin{lemma}\label{lem:Z4ss}
Let $E$ be the supersingular elliptic curve over $k$. Every faithful $\Z/4$-action on
$E$ is conjugate in $\operatorname{Aut}(E)$ to one fixing $O$, and the six elements of
order $4$ of $\operatorname{Aut}(E,O)\cong\mathrm{SL}_2(\mathbb F_3)$ form one conjugacy
class; so up to isomorphism of $G$-curves (and replacing the generator by its inverse)
$(E,\Z/4)$ is $(E_0,\text{Witt action})$. There is no faithful $\Z/8$-action on $E$
(Lemma~2.5(iii)). \stag{proved}
\end{lemma}

\begin{proof}
Write $\varphi=t_P\circ\alpha$ with $\alpha\in\operatorname{Aut}(E,O)$. If $\alpha=1$,
$\operatorname{ord}\varphi=\operatorname{ord}P$ is odd ($E(k)[2]=0$); if $\alpha=-1$,
$\varphi^2=t_{P+\alpha P}=1$. So order $4$ forces $\operatorname{ord}\alpha=4$, and
$t_Q\varphi t_{-Q}=t_{(1-\alpha)Q}\alpha$ with $1-\alpha$ a nonzero isogeny, surjective
on $E(k)$; choose $Q$ with $(1-\alpha)Q=-P$. The conjugacy classes of
$\mathrm{SL}_2(\mathbb F_3)$ are $\{1\}$, $\{-1\}$, two of order $3$, two of order $6$
and one of order $4$, of size $6$.
\end{proof}

\emph{Enumeration} (\texttt{rh\_enum.py}, output \texttt{rh\_enum.out}) \stag{computed}.
Imposing only the necessary conditions above (integrality and non-negativity of genus and
$p$-rank at every level of the tower $C\to C/\langle g^{2^{n-1}}\rangle\to\dots\to C/G$,
Hasse--Arf, $u_{i+1}\ge2u_i$), the possible $(g(C),g(C/G),\sigma(C/G),\sigma(C))$ for
faithful $\Z/8$-curves with $g(C)\le10$ are:
\begin{center}\small
\begin{tabular}{@{}llll@{}}
\toprule
$g(C)$ & $C/G$ & ramification & $\sigma(C)$\\
\midrule
$1$ & ordinary elliptic & free action (translation by an $8$-torsion point) & $1$\\
$5$ & ordinary elliptic & one orbit with inertia $\Z/2$ & $5$\\
$7$ & $\mathbb P^1$ & one totally ramified point, lower breaks $(1,3,11)$ & $0$\\
$9$ & $\mathbb P^1$ & one totally ramified point, lower breaks $(1,3,15)$ & $0$\\
$9$ & genus $1$ or $2$ & various (one or two orbits, or free) & $5,7,9$ or $1$\\
\bottomrule
\end{tabular}
\end{center}
Consequences: a faithful $\Z/8$-curve of genus $\le6$ is an ordinary elliptic curve
with translation action or has genus $5$ and $p$-rank $5$; no faithful $\Z/8$-curve has
genus $2,3,4,6,8$; and a $p$-rank-$0$ one of genus $\le10$ has genus $7$ or $9$, with
$C/G=\mathbb P^1$, $\Z/4$-quotient a supersingular elliptic curve and $\Z/2$-quotient
$\mathbb P^1$. For $\Z/4$ and $g(C)\le6$: the $p$-rank-$0$ entries have a single totally
ramified point with lower breaks $(1,2j+1)$ (genus $j$) or $C/G$ of genus $1$; the only
other genus-$1$ entry is an ordinary elliptic curve with translation by a $4$-torsion
point. A genus-$2$ faithful $\Z/4$-curve has $p$-rank $0$, quotient $\mathbb P^1$ and
lower breaks $(1,5)$.

\section{The level-\texorpdfstring{$2$}{2} set \texorpdfstring{$\mathfrak F$}{F} and the
isogeny mechanism}\label{app:isog}

Recall $\mathfrak F=\{f\in K/\wp K:\ed_k((s_0,s_1)_K\otimes K(v_f))=1\}$
(Theorem~\ref{thm:caseA}(c)). Call a descent of $(s_0,s_1)_K\otimes K(v_f)$ ``of Case A''
if its descent field lies inside $K$.

\begin{proposition}[level-$2$ Case A is explicit]\label{prop:FA}
$f$ admits a level-$2$ Case-A descent, i.e.\ $(s_0,s_1)\equiv\beta(t)$ over $K(v_f)$
with $t\in K$ and $\beta\in W_2(k(T))$, iff
\[
f\equiv s_1+x_0(x_0+1)(x_0+s_0)+\beta_1(t)\pmod{\wp K}
\]
for some $x_0,t\in K$ and $\beta_0,\beta_1\in k(T)$ with $\beta_0(t)=s_0+x_0^2+x_0$.
This set $\mathfrak F_A$ contains $\{s_1+x_0(x_0+1)(x_0+s_0)+b(s_0+x_0^2+x_0):x_0\in K,\
b\in k(T)\}$ (the case $\beta_0=T$). \stag{proved}
\end{proposition}

\begin{proof}
``$\Leftarrow$'': over $K'=K(v_f)$ put $x_1:=v_f+h$ where $f=s_1+x_0(x_0+1)(x_0+s_0)+
\beta_1(t)+\wp h$; by Appendix~\ref{app:witt},
$(s_0,s_1)+\wp(x_0,x_1)=(\beta_0(t),\beta_1(t))$. ``$\Rightarrow$'': the descent field is
$k(t)$ (Theorem~\ref{thm:caseA}(a)); $(s_0,s_1)$ and $\beta(t)$ agree over $K'$, so
$(s_0,s_1)\equiv\beta(t)+\varepsilon(0,f)$ over $K$ with $\varepsilon=1$ since
$\ed_k((s_0,s_1)_K)=2$; read off the coordinates of $(s_0,s_1)+\wp(x_0,x_1)=
(\beta_0(t),\beta_1(t)+f)$.
\end{proof}

\begin{proposition}[the $E_0$-part of $\mathfrak F$]\label{prop:FE0}
$\{f:E'(K(v_f))\ne\{\infty'\}\}=\{s_1+u(u+1)(u+s_0):u\in K\}+\wp K\subset\mathfrak F_A$.
Equivalently, the quadratic twist $E'^{(f)}\colon w^2+w=u(u+1)(u+s_0)+s_1+f$ has an
affine $K$-point iff $f$ is in this family, and never one with $u\notin K$. \stag{proved}
\end{proposition}

\begin{proof}
As in Corollary~\ref{cor:noquad}, every $Q\in E'(K')$ has $\sigma Q=-Q$, so $u(Q)\in K$;
then $w\in K'\smallsetminus K$ with $w^2+w\in K$ forces $w=a+v_f$, $a\in K$, and
$f\equiv s_1+u(u+1)(u+s_0)$. Conversely such $f$ gives the point $(u,v_f+a)$.
\end{proof}

\begin{proposition}[$f\in k(s_0)$]\label{prop:ks0}
For $e\in k(s_0)$ with $e\notin\{0,s_0\}+\wp k(s_0)$: $\ed_k(\tau_2\otimes K_2(v_e))=2$.
Hence $k(s_0)\cap\mathfrak F=\{s_0\}+\wp K$. \stag{proved}
\end{proposition}

\begin{proof}
The total space of $\tau_2\otimes K_2(v_e)$ is $k(z_0,v_e)(z_1)=M(z_1)$ with
$M=k(z_0,v_e)$ fixed by $g^2\colon(z_0,z_1)\mapsto(z_0,z_1+1)$; apply
Proposition~\ref{prop:fft} for $\Z/4$. (No genus hypothesis on the curve
$v^2+v=e(z_0^2+z_0)$ is needed; for $e\in k(s_0)$ the twist $E'^{(e)}$ is isomorphic to
$E'$ via $s_1\mapsto s_1+e$, which shows directly $E'^{(e)}(K_2)=0$, consistent with
Proposition~\ref{prop:FE0}.)
\end{proof}

\begin{proposition}[corestriction exclusion of constant descents]\label{prop:cores}
Let $M_0\subset K$ be a subfield algebraically closed in $K$, $f'\in M_0$ with
$f\equiv f'$ mod $\wp K$, $M=M_0(v_{f'})$, so $K'=K\otimes_{M_0}M$. If $s_0\notin
M_0+\wp K$ (e.g.\ $M_0\subset k(s_1,s_2)$), then no descent field of
$(s_0,s_1)_K\otimes K'$, in particular none of $\tau_\gen\otimes K'$, is contained in
$M$. \stag{proved}
\end{proposition}

\begin{proof}
Let $(s_0,s_1)\equiv\beta$ over $K'$ with $\beta\in W_2(M)/\wp$. Corestriction is
compatible with base change along $M_0\to K$, so $\cores_{K'/K}(\beta)=\gamma$ comes from
$W_2(M_0)/\wp$, while $\cores_{K'/K}\res(s_0,s_1)=2(s_0,s_1)=VF(s_0,s_1)=(0,s_0^2)\equiv
(0,s_0)$. So $(0,s_0)\equiv(\gamma_0,\gamma_1)$ over $K$; $\gamma_0\in\wp K\cap M_0=
\wp M_0$, so after modifying $\gamma$ by $\wp(h,0)$, $\gamma_0=0$ and
$s_0+\gamma_1\in\wp K$ by (5), contradicting $s_0\notin M_0+\wp K$.
\end{proof}

This does not exclude a general Case-B descent: for a $\sigma$-stable descent field it
only says that $(0,s_0)=\operatorname{Ind}(s_0)$, which already has $\ed=1$, descends to
it.

\emph{The isogeny mechanism.} By Theorem~\ref{thm:struct}(a) and Appendix~\ref{app:RH},
a witness curve of genus $\le6$ is an ordinary elliptic curve $C$ with $G$ acting by
translation by a point $P$ of order $8$, so $C'=C/\langle P\rangle$ is ordinary elliptic,
$C\to C'$ is the \'etale $\Z/8$-isogeny (the dual of $V^3$, i.e.\ $C=C'^{(8)}$), and $C'$
is dominated by $D_{f,a}$ for general $a$: the Jacobians of the curves $D_{f,a}$ must have
a \emph{fixed} ordinary elliptic factor, which requires $f(a,s_2)$ to have at least two
poles on the $s_2$-line. The class of $\tau_\gen\otimes K'$ is then $\delta(Q)$ for the
connecting map $\delta\colon C'(K')\to H^1(K',\Z/8)$ of $0\to\Z/8\to C\to C'\to0$ and
the non-constant point $Q\in C'(K')$ given by $k(C')\subset K'$. For $n=1$ and
$C'\colon y^2+xy=x^3+a_6$ (ordinary, $j=1/a_6$), the substitution $y=xz$ gives
$z^2+z=x+a_6/x^2$, so the \'etale $\Z/2$-isogeny torsor at $(x,y)$ has Artin--Schreier
class $x+a_6/x^2$; the $\Z/8$ version is the Witt vector of the tower
$C'^{(8)}\to C'^{(4)}\to C'^{(2)}\to C'$, not computed here. Since $\delta$ is a
homomorphism and $s$ is $\sigma$-invariant, $\delta(Q+\sigma Q)=2s\equiv(0,0,s_0)$ and
$\delta(Q-\sigma Q)=0$, i.e.\ $Q-\sigma Q$ lifts to $C(K')$. We have not derived a
contradiction from these constraints. \stag{open}

\section{The cubic slot on \texorpdfstring{$\bar U_3$}{U\_3}: Witt trace, chain
conditions, searches}\label{app:cubic}

\emph{Witt trace.} For a cubic $u^3+Au^2+Bu+C$ with $A=1+s_0+\alpha^2$, $B=s_0+\alpha$,
$C=s_1+\beta^2+\beta$ and $t=s_0+u^2+u$, the trace $[t]+[t']+[t'']$ has coordinates
$(e_1,e_2,e_1e_3+e_1^2e_2)$ in the elementary symmetric functions of $t,t',t''$
(from (3); the ghost components of $\sum[t_i]$ are the power sums $\sum t_i^{2^m}$).
Moreover $3x=(x_0,x_1+x_0^2,x_2+x_1^2+x_0^2x_1)$. The scripts \texttt{cores\_cubic*.py}
compute $d:=3s-\operatorname{Tr}[t]$: $d_0=\wp(A)$; after subtracting $\wp[A]$ the
coordinate $1$ is $\wp(\alpha^4+\alpha^3+\alpha s_0+\beta+s_0^2+s_0)$, automatically in
$\wp K$ (the $n=2$ statement of Theorem~\ref{thm:A} seen through $\cores$); after
subtracting $\wp V[y_1]$ the last coordinate is $s_2+R(s_0,s_1,\alpha,\beta)$ with $R$ a
polynomial of $67$ terms, and it agrees with $\operatorname{Tr}_{K(u)/K}(e_Q)$ modulo
$\wp K$. Reducing squares ($x^2\equiv x$), $R\equiv R_{\rm red}$ with
{\small\begin{align*}
R_{\rm red}={}&\alpha^{11}+s_0\alpha^{10}+s_0\alpha^9+\alpha^8\beta+s_1\alpha^8+s_0\alpha^7
+\alpha^6\beta+s_1\alpha^6+s_0\alpha^6+\alpha^5\beta^2+\alpha^5\beta+\alpha^5+s_1\alpha^5+s_0^2\alpha^5\\
&+s_0\alpha^4\beta^2+\alpha^4\beta+s_0\alpha^4\beta+s_0s_1\alpha^4+\alpha^3\beta
+\alpha^3+s_0\alpha^3+s_0^3\alpha^3+s_0^4\alpha^3+s_0^2\alpha^2\beta+s_0\alpha^2+s_0^2s_1\alpha^2\\
&+s_0^3\alpha^2+s_0^5\alpha^2+\alpha\beta^2+s_0\alpha\beta^2+s_0^2\alpha\beta^2
+\alpha\beta+s_0^2\alpha\beta+\alpha+s_1\alpha+s_0\alpha+s_0s_1\alpha+s_0^2s_1\alpha+s_0^4\alpha\\
&+s_0^5\alpha+\beta^3+s_1\beta^2+s_0\beta^2+s_0^3\beta^2+\beta+s_1\beta
+s_0^3\beta+s_0^4\beta+s_1+s_0s_1+s_0^2s_1+s_0^3s_1+s_0^4s_1 .
\end{align*}}
The twisted $G$-action on $W_3$ acts on lines by $(\alpha,\beta)\mapsto(\alpha+1,
\alpha+\beta+1)$, $(\alpha,\beta+1)$ and preserves \eqref{eq:N1} (checked symbolically).
\stag{computed}

\emph{Chain conditions at $s_2=\infty$.} Let $\hat K=K_2((1/s_2))$ and
$e=\sum_{i\le N}c_is_2^i\in\hat K$, $c_i\in K_2$. Then $e\in\wp\hat K$ iff $c_0\in\wp K_2$
and, for every odd $i_0\ge1$, with $2^Ji_0$ the last index of the chain
$i_0,2i_0,4i_0,\dots$ carrying a nonzero coefficient,
\[
(\mathrm{Ch}_{i_0})\qquad\sum_{j=0}^{J}c_{i_02^j}^{\,2^{J-j}}=0\quad\text{in }K_2 .
\]
(The negative part is always $\wp$-exact; the solution $y=\sum b_is_2^i$ of
$y^2+y=\sum_{i>0}c_is_2^i$ has $b_{i_0}=c_{i_0}$, $b_{2m}=c_{2m}+b_m^2$, and must be a
finite sum.) In particular a top coefficient $c_N\ne0$ with $N$ odd excludes
$e\in\wp K$. Proposition~\ref{prop:D}(b) uses: for $\alpha=a_0+a_1s_2$,
$\beta=b_0+b_1s_2$ (plus arbitrary tails of negative order), $c_{11}=a_1^{11}$ and
$c_{13}=c_{14}=c_{15}=0$; with $a_1=0$, $c_3=b_1^3$ and $c_6=c_{12}=0$
(\texttt{deg1\_family.py}). For pole orders $\ge2$ the chain equations are nonlinear in
the coefficients and Laurent tails enter; no general argument is known to us.

\emph{Searches} (\texttt{f2poly.py}, \texttt{search\_lines.py},
\texttt{search\_targeted.py}) \stag{computed}. For $\alpha,\beta\in\mathbb F_2[s_0,s_1,s_2]$,
membership $s_2+R\in\wp K$ is equivalent to membership in $\wp(k[s_0,s_1,s_2])$ and is
decided exactly by the monomial-chain algorithm (top monomial must be a square; subtract
$\wp$ of its root; iterate). None of the following families contains a line satisfying
\eqref{eq:N1}: $\deg\alpha\le1$, $\deg\beta\le2$ ($16384$ lines); $\alpha=0$,
$\deg\beta\le3$ ($2^{20}$ lines); $\alpha\in\mathbb F_2[s_0,s_1]$ of degree $\le2$ with
$\beta$ of $s_2$-degree $\le2$ and $(s_0,s_1)$-degree $\le1$ ($2^{15}$); $\alpha$ of
$s_2$-degree $\le1$ and $(s_0,s_1)$-degree $\le1$ with $\beta$ of $s_2$-degree $\le2$
resp.\ $\le3$ ($2^{15}$ resp.\ $2^{18}$); $\alpha,\beta$ both of total degree $\le2$
($2^{20}$). A finite-field specialisation test (\texttt{ff\_test.py}: specialise
$(s_0,s_1,s_2)\in\mathbb F_8^3$, take the roots $u\in\mathbb F_{512}$ of the cubic and
compute the absolute trace of $e_Q(u)$) is available as an independent check of the full
condition $e_Q\in\wp K(u)$ on any candidate; it has not been needed.

\section{Scripts}\label{app:scripts}
All scripts are Python~3 and live in \texttt{notes/scripts/} of the repository:
\texttt{witt\_check.py} (ghost-component verification of (3)--(4), prints $\Phi_2$),
\texttt{wittlib.py}, \texttt{ed2\_witt.py} (negation, $s\pm\wp x$),
\texttt{genus\_U3.py} (Artin--Schreier reduction on $E_0$, $g(\bar U_3)=7$),
\texttt{rh\_enum.py} (ramification enumeration, output \texttt{rh\_enum.out}),
\texttt{cores\_cubic.py}, \texttt{cores\_cubic2.py}, \texttt{cores\_cubic3.py},
\texttt{resolvent.py} (Witt trace, $R$, invariance under the twisted action),
\texttt{build\_R.py}, \texttt{f2poly.py}, \texttt{search\_lines.py},
\texttt{search\_targeted.py}, \texttt{deg1\_family.py}, \texttt{ff\_test.py}.


\begin{thebibliography}{99}
\setcounter{enumiv}{16}
\bibitem{Ar} J.~Kr.~Arason, \emph{Cohomologische Invarianten quadratischer Formen},
J. Algebra 36 (1975), 448--491.
\bibitem{BF} G.~Berhuy and G.~Favi, \emph{Essential dimension: a functorial point of
view (after A.~Merkurjev)}, Doc. Math. 8 (2003), 279--330. (= [1])
\bibitem{BR} J.~Buhler and Z.~Reichstein, \emph{On the essential dimension of a finite
group}, Compositio Math. 106 (1997), 159--179. (= [2])
\bibitem{Cr} R.~Crew, \emph{\'Etale $p$-covers in characteristic $p$}, Compositio Math.
52 (1984), 31--45.
\bibitem{DR} A.~Duncan and Z.~Reichstein, \emph{Versality of algebraic group actions
and rational points on twisted varieties}, J. Algebraic Geom. 24 (2015), 499--530.
(= [3])
\bibitem{FJ} M.~D.~Fried and M.~Jarden, \emph{Field arithmetic}, 3rd ed., Springer,
2008.
\bibitem{Ga} M.~Garuti, \emph{Linear systems attached to cyclic inertia}, in:
Arithmetic fundamental groups and noncommutative algebra, Proc. Sympos. Pure Math. 70,
AMS, 2002, pp. 377--386.
\bibitem{KM} N.~Karpenko and A.~Merkurjev, \emph{Essential dimension of finite
$p$-groups}, Invent. Math. 172 (2008), 491--508.
\bibitem{LT} S.~Lang and J.~Tate, \emph{Principal homogeneous spaces over abelian
varieties}, Amer. J. Math. 80 (1958), 659--684.
\bibitem{Le07} A.~Ledet, \emph{Finite groups of essential dimension one}, J. Algebra
311 (2007), 31--37. (= [6])
\bibitem{Me} A.~S.~Merkurjev, \emph{Essential dimension: a survey}, Transform. Groups
18 (2013), 415--481. (= [9])
\bibitem{MR} A.~Meyer and Z.~Reichstein, \emph{Some consequences of the
Karpenko--Merkurjev theorem}, Doc. Math. Extra Vol. (2010), 445--457.
\bibitem{RS} Z.~Reichstein and F.~Scavia, \emph{The behavior of essential dimension
under specialization}, \'Epijournal G\'eom. Alg\'ebrique 6 (2022), Art. 21. (= [12])
\bibitem{Se} J.-P.~Serre, \emph{Galois cohomology}, Springer, 1997.
\bibitem{Th} L.~Thomas, \emph{Ramification groups in Artin--Schreier--Witt extensions},
J. Th\'eor. Nombres Bordeaux 17 (2005), 689--720.
\end{thebibliography}

\end{document}
